\begin{register}{H}{MCPWM\_CLK\_CFG\_REG}{0x0000}\label{regdesc:MCPWMCLKCFGREG}
\regfield{(reserved)}{24}{8}{000000000000000000000000}%
\regfield{MCPWM\_CLK\_PRESCALE}{8}{0}{{0x0}}%
\reglabel{Reset}\regnewline%
\begin{regdesc}\begin{reglist}
\label{fielddesc:MCPWMCLKPRESCALE}\item [MCPWM\_CLK\_PRESCALE] 配置时钟预分频系数，使得 PWM\_CLK 的周期为 6.25 ns * (PWM\_CLK\_PRESCALE + 1)。 (R/W)
\end{reglist}\end{regdesc}
\end{register}


\begin{register}{H}{MCPWM\_TIMER\regindex{n}\_CFG0\_REG (\regindex{n}: 0-2)}{0x0004+0x10*\regindex{n}}\label{regdesc:MCPWMTIMERNCFG0REG}
\regfield{(reserved)}{6}{26}{000000}%
\regfield{MCPWM\_TIMER\regindex{n}\_PERIOD\_UPMETHOD}{2}{24}{{0}}%
\regfield{MCPWM\_TIMER\regindex{n}\_PERIOD}{16}{8}{{0xff}}%
\regfield{MCPWM\_TIMER\regindex{n}\_PRESCALE}{8}{0}{{0x0}}%
\reglabel{Reset}\regnewline%
\begin{regdesc}\begin{reglist}
\label{fielddesc:MCPWMTIMERNPRESCALE}\item [MCPWM\_TIMER\regindex{n}\_PRESCALE] 配置定时器\regindex{n} 的预分频系数，使得 PT0\_CLK 的周期 = PWM\_CLK 的周期 * (PWM\_TIMER\regindex{n}\_PRESCALE + 1)。\\
(R/W)
\label{fielddesc:MCPWMTIMERNPERIOD}\item [MCPWM\_TIMER\regindex{n}\_PERIOD] 配置定时器\regindex{n} 的影子周期。\\
(R/W)
\label{fielddesc:MCPWMTIMERNPERIODUPMETHOD}\item [MCPWM\_TIMER\regindex{n}\_PERIOD\_UPMETHOD] 配置 PWM 定时器\regindex{n} 周期有效寄存器的更新方式。\\
0：立即更新\\
1：发生 TEZ 事件时更新\\
2：同步时更新\\
3：发生 TEZ 事件或同步时更新\\
本文档中，TEZ 指定时器为 0 时的事件。\\ (R/W)
\end{reglist}\end{regdesc}
\end{register}


\begin{register}{H}{MCPWM\_TIMER\regindex{n}\_CFG1\_REG(\regindex{n}: 0-2)}
{0x0008+0x10*\regindex{n}}\label{regdesc:MCPWMTIMERNCFG1REG}
\regfield{(reserved)}{27}{5}{000000000000000000000000000}%
\regfield{MCPWM\_TIMER\regindex{n}\_MOD}{2}{3}{{0x0}}%
\regfield{MCPWM\_TIMER\regindex{n}\_START}{3}{0}{{0x0}}%
\reglabel{Reset}\regnewline%
\begin{regdesc}\begin{reglist}
\label{fielddesc:MCPWMTIMERNSTART}\item [MCPWM\_TIMER\regindex{n}\_START] 控制 PWM 定时器的开启与关闭。\\
0：如果开启，在 TEZ 事件发生时停止\\
1：如果开启，在 TEP 事件发生时停止\\
2：开启\\
3：开启，并在下一个 TEZ 事件发生时停止\\
4：开启，并在下一个 TEP 事件发生时停止\\
其他值：无效值，不产生影响\\
本文档中，TEP 指定时器为周期值时发生的事件。\\ (R/W/SC)
\label{fielddesc:MCPWMTIMERNMOD}\item [MCPWM\_TIMER\regindex{n}\_MOD] 配置 PWM 定时器 0 的工作模式。\\
0：暂停\\
1：递增模式\\
2：递减模式\\
3：递增递减循环模式\\
(R/W)
\end{reglist}\end{regdesc}
\end{register}


\begin{register}{H}{MCPWM\_TIMER\regindex{n}\_SYNC\_REG
(\regindex{n}: 0-2)}
{0x000C+0x10*\regindex{n}}\label{regdesc:MCPWMTIMERNSYNCREG}
\regfield{(reserved)}{11}{21}{00000000000}%
\regfield{MCPWM\_TIMER\regindex{n}\_PHASE\_DIRECTION}{1}{20}{{0}}%
\regfield{MCPWM\_TIMER\regindex{n}\_PHASE}{16}{4}{{0}}%
\regfield{MCPWM\_TIMER\regindex{n}\_SYNCO\_SEL}{2}{2}{{0}}%
\regfield{MCPWM\_TIMER\regindex{n}\_SYNC\_SW}{1}{1}{{0}}%
\regfield{MCPWM\_TIMER\regindex{n}\_SYNCI\_EN}{1}{0}{{0}}%
\reglabel{Reset}\regnewline%
\begin{regdesc}\begin{reglist}
\label{fielddesc:MCPWMTIMERNSYNCIEN}\item [MCPWM\_TIMER\regindex{n}\_SYNCI\_EN] 配置是否开启同步输入事件发生时的定时器\regindex{n} 相位重载。
\\0：关闭
\\1：开启
\\(R/W)
\label{fielddesc:MCPWMTIMERNSYNCSW}\item [MCPWM\_TIMER\regindex{n}\_SYNC\_SW] 配置是否触发软件同步。\\
0：无效\\
1：触发软件同步\\
(R/W)
\label{fielddesc:MCPWMTIMERNSYNCOSEL}\item [MCPWM\_TIMER\regindex{n}\_SYNCO\_SEL] 选择 PWM 定时器\regindex{n} 的同步输出来源。\\
0：同步。取反 \hyperref[fielddesc:MCPWMTIMERNSYNCSW]{MCPWM\_TIMER\regindex{n}\_SYNC\_SW} 位时始终生成同步输出。\\
1：TEZ\\
2：TEP\\
3：无效\\
(R/W)
\label{fielddesc:MCPWMTIMERNPHASE}\item [MCPWM\_TIMER\regindex{n}\_PHASE] 同步事件中定时器\regindex{n} 重载的相位。 \\(R/W)
\label{fielddesc:MCPWMTIMERNPHASEDIRECTION}\item [MCPWM\_TIMER\regindex{n}\_PHASE\_DIRECTION] 当定时器\regindex{n} 为递增-递减循环模式时，配置该定时器方向。\\
0：递增\\
1：递减\\
(R/W)
\end{reglist}\end{regdesc}
\end{register}


\begin{register}{H}{MCPWM\_TIMER\regindex{n}\_STATUS\_REG(\regindex{n}: 0-2)}
{0x0010+0x10*\regindex{n}}\label{regdesc:MCPWMTIMERNSTATUSREG}
\regfield{(reserved)}{15}{17}{000000000000000}%
\regfield{MCPWM\_TIMER\regindex{n}\_DIRECTION}{1}{16}{{0}}%
\regfield{MCPWM\_TIMER\regindex{n}\_VALUE}{16}{0}{{0}}%
\reglabel{Reset}\regnewline%
\begin{regdesc}\begin{reglist}
\label{fielddesc:MCPWMTIMERNVALUE}\item [MCPWM\_TIMER\regindex{n}\_VALUE] 表示当前 PWM 定时器\regindex{n} 计数器的值。 (RO)
\label{fielddesc:MCPWMTIMERNDIRECTION}\item [MCPWM\_TIMER\regindex{n}\_DIRECTION] 配置当前 PWM 定时器\regindex{n} 的计数器模式。\\
0：递增模式\\
1：递减模式\\
(RO)
\end{reglist}\end{regdesc}
\end{register}


\begin{register}{H}{MCPWM\_TIMER\_SYNCI\_CFG\_REG}{0x0034}\label{regdesc:MCPWMTIMERSYNCICFGREG}
\regfield{(reserved)}{20}{12}{00000000000000000000}%
\regfield{MCPWM\_EXTERNAL\_SYNCI2\_INVERT}{1}{11}{{0}}%
\regfield{MCPWM\_EXTERNAL\_SYNCI1\_INVERT}{1}{10}{{0}}%
\regfield{MCPWM\_EXTERNAL\_SYNCI0\_INVERT}{1}{9}{{0}}%
\regfield{MCPWM\_TIMER2\_SYNCISEL}{3}{6}{{0}}%
\regfield{MCPWM\_TIMER1\_SYNCISEL}{3}{3}{{0}}%
\regfield{MCPWM\_TIMER0\_SYNCISEL}{3}{0}{{0}}%
\reglabel{Reset}\regnewline%
\begin{regdesc}\begin{reglist}
\label{fielddesc:MCPWMTIMERNSYNCISEL}\item [MCPWM\_TIMER\regindex{n}\_SYNCISEL] 选择 PWM 定时器\regindex{n} 的同步输入来源。 \\
1：PWM 定时器 0 同步输出\\
2：PWM 定时器 1 同步输出\\
3：PWM 定时器 2 同步输出\\
4：来自 GPIO 矩阵的 SYNC0\\
5：来自 GPIO 矩阵的 SYNC1\\
6：来自 GPIO 矩阵的 SYNC2\\
其他值：未选择任何同步输入\\
(R/W)
\label{fielddesc:MCPWMEXTERNALSYNCININVERT}\item [MCPWM\_EXTERNAL\_SYNCI\regindex{n}\_INVERT] 配置是否将来自 GPIO 矩阵的 SYNC\regindex{n} 反相。\\
0：不反相\\
1：反相 \\
(R/W)
\end{reglist}\end{regdesc}
\end{register}


\begin{register}{H}{MCPWM\_OPERATOR\_TIMERSEL\_REG}{0x0038}\label{regdesc:MCPWMOPERATORTIMERSELREG}
\regfield{(reserved)}{26}{6}{00000000000000000000000000}%
\regfield{MCPWM\_OPERATOR2\_TIMERSEL}{2}{4}{{0}}%
\regfield{MCPWM\_OPERATOR1\_TIMERSEL}{2}{2}{{0}}%
\regfield{MCPWM\_OPERATOR0\_TIMERSEL}{2}{0}{{0}}%
\reglabel{Reset}\regnewline%
\begin{regdesc}\begin{reglist}
\label{fielddesc:MCPWMOPERATORNTIMERSEL}\item [MCPWM\_OPERATOR\regindex{n}\_TIMERSEL] 选择 PWM 操作器\regindex{n} 的定时参考来源。\\
0：定时器 0\\
1：定时器 1\\
2：定时器 2\\
3：无效，将选择定时器 2\\
(R/W)
\end{reglist}\end{regdesc}
\end{register}


\begin{register}{H}{MCPWM\_GEN\regindex{n}\_STMP\_CFG\_REG (\regindex{n}: 0-2)}{0x003C+0x38*\regindex{n}}\label{regdesc:MCPWMGENNSTMPCFGREG}
\regfield{(reserved)}{22}{10}{0000000000000000000000}%
\regfield{MCPWM\_CMPR\regindex{n}\_B\_SHDW\_FULL}{1}{9}{{0}}%
\regfield{MCPWM\_CMPR\regindex{n}\_A\_SHDW\_FULL}{1}{8}{{0}}%
\regfield{MCPWM\_CMPR\regindex{n}\_B\_UPMETHOD}{4}{4}{{0}}%
\regfield{MCPWM\_CMPR\regindex{n}\_A\_UPMETHOD}{4}{0}{{0}}%
\reglabel{Reset}\regnewline%
\begin{regdesc}\begin{reglist}
\label{fielddesc:MCPWMCMPRNAUPMETHOD}\item [MCPWM\_CMPR\regindex{n}\_A\_UPMETHOD] 配置 PWM 生成器\regindex{n} 时间戳寄存器 A 的有效寄存器的更新方式。\\
所有位值为 0：立即更新\\
bit0 为 1：发生 TEZ 事件时更新\\
bit1 为 1：发生 TEP 事件时更新\\
bit2 为 1：发生同步时间时更新\\
bit3 为 1：关闭更新\\
(R/W)
\label{fielddesc:MCPWMCMPRNBUPMETHOD}\item [MCPWM\_CMPR\regindex{n}\_B\_UPMETHOD] 配置 PWM 生成器\regindex{n} 时间戳寄存器 B 的有效寄存器的更新方式。\\
所有位值为 0：立即更新\\
bit0 为 1：发生 TEZ 事件时更新\\
bit1 为 1：发生 TEP 事件时更新\\
bit2 为 1：发生同步时间时更新\\
bit3 为 1：关闭更新\\
(R/W)
\label{fielddesc:MCPWMCMPRNASHDWFULL}\item [MCPWM\_CMPR\regindex{n}\_A\_SHDW\_FULL] 配置是否在某个特定时间将影子寄存器中的值写入对应的有效寄存器中。由硬件置 1 或复位。\\
0：将影子寄存器中最新的值写入有效寄存器 A 中\\
1：将值写入 PWM 生成器\regindex{n} 时间戳寄存器 A 的影子寄存器中，等待传输给有效寄存器 A。\\
(R/SC/WTC)
\label{fielddesc:MCPWMCMPRNBSHDWFULL}\item [MCPWM\_CMPR\regindex{n}\_B\_SHDW\_FULL] 配置是否在某个特定时间将影子寄存器中的值写入对应的有效寄存器中。由硬件置 1 或复位。\\
0：将影子寄存器中最新的值写入有效寄存器 B 中\\
1：将值写入 PWM 生成器\regindex{n} 时间戳寄存器 B 的影子寄存器中，等待传输给有效寄存器 B。\\
(R/SC/WTC)
\end{reglist}\end{regdesc}
\end{register}


\begin{register}{H}{MCPWM\_GEN\regindex{n}\_TSTMP\_A\_REG(\regindex{n}: 0-2)}
{0x0040+0x38*\regindex{n}}\label{regdesc:MCPWMGENNTSTMPAREG}
\regfield{(reserved)}{16}{16}{0000000000000000}%
\regfield{MCPWM\_CMPR\regindex{n}\_A}{16}{0}{{0}}%
\reglabel{Reset}\regnewline%
\begin{regdesc}\begin{reglist}
\label{fielddesc:MCPWMCMPRNA}\item [MCPWM\_CMPR\regindex{n}\_A] 配置 PWM 生成器\regindex{n} 时间戳寄存器 A 的影子寄存器。\\(R/W)
\end{reglist}\end{regdesc}
\end{register}


\begin{register}{H}{MCPWM\_GEN\regindex{n}\_TSTMP\_B\_REG(\regindex{n}: 0-2)}
{0x0044+0x38*\regindex{n}}\label{regdesc:MCPWMGENNTSTMPBREG}
\regfield{(reserved)}{16}{16}{0000000000000000}%
\regfield{MCPWM\_CMPR\regindex{n}\_B}{16}{0}{{0}}%
\reglabel{Reset}\regnewline%
\begin{regdesc}\begin{reglist}
\label{fielddesc:MCPWMCMPRNB}\item [MCPWM\_CMPR\regindex{n}\_B] 配置 PWM 生成器\regindex{n} 时间戳寄存器 B 的影子寄存器。 \\(R/W)
\end{reglist}\end{regdesc}
\end{register}


\begin{register}{H}{MCPWM\_GEN\regindex{n}\_CFG0\_REG(\regindex{n}: 0-2)}
{0x0048+0x38*\regindex{n}}\label{regdesc:MCPWMGENNCFG0REG}
\regfield{(reserved)}{22}{10}{0000000000000000000000}%
\regfield{MCPWM\_GEN\regindex{n}\_T1\_SEL}{3}{7}{{0}}%
\regfield{MCPWM\_GEN\regindex{n}\_T0\_SEL}{3}{4}{{0}}%
\regfield{MCPWM\_GEN\regindex{n}\_CFG\_UPMETHOD}{4}{0}{{0}}%
\reglabel{Reset}\regnewline%
\begin{regdesc}\begin{reglist}
\label{fielddesc:MCPWMGENNCFGUPMETHOD}\item [MCPWM\_GEN\regindex{n}\_CFG\_UPMETHOD] 配置 PWM 生成器\regindex{n} 有效配置寄存器的更新方式。\\
所有 bit 值为 0：立即更新\\
bit0 为 1：发生 TEZ 事件时更新\\
bit1 为 1：发生 TEP 事件时更新\\
bit2 为 1：发生同 步时间时更新\\
bit3 为 1：关闭更新\\
(R/W)
\label{fielddesc:MCPWMGENNT0SEL}\item [MCPWM\_GEN\regindex{n}\_T0\_SEL] 选择 PWM 生成器\regindex{n} event\_t0 的信号源，立即生效。\\
0：fault\_event0\\
1：fault\_event1\\
2：fault\_event2\\
3：sync\_taken\\
4：无\\
(R/W)
\label{fielddesc:MCPWMGENNT1SEL}\item [MCPWM\_GEN\regindex{n}\_T1\_SEL] 选择 PWM 生成器\regindex{n} event\_t1 的信号源，立即生效。\\
0：fault\_event0\\
1：fault\_event1\\
2：fault\_event2\\
3：sync\_taken\\
4：无\\
(R/W)
\end{reglist}\end{regdesc}
\end{register}

\begin{register}{H}{MCPWM\_GEN\regindex{n}\_FORCE\_REG(\regindex{n}: 0-2)}
{0x004C+0x38*\regindex{n}}\label{regdesc:MCPWMGENNFORCEREG}
\regfield{(reserved)}{16}{16}{0000000000000000}%
\regfield{MCPWM\_GEN\regindex{n}\_B\_NCIFORCE\_MODE}{2}{14}{{0}}%
\regfield{MCPWM\_GEN\regindex{n}\_B\_NCIFORCE}{1}{13}{{0}}%
\regfield{MCPWM\_GEN\regindex{n}\_A\_NCIFORCE\_MODE}{2}{11}{{0}}%
\regfield{MCPWM\_GEN\regindex{n}\_A\_NCIFORCE}{1}{10}{{0}}%
\regfield{MCPWM\_GEN\regindex{n}\_B\_CNTUFORCE\_MODE}{2}{8}{{0}}%
\regfield{MCPWM\_GEN\regindex{n}\_A\_CNTUFORCE\_MODE}{2}{6}{{0}}%
\regfield{MCPWM\_GEN\regindex{n}\_CNTUFORCE\_UPMETHOD}{6}{0}{{0x20}}%
\reglabel{Reset}\regnewline%
\begin{regdesc}\begin{reglist}
\label{fielddesc:MCPWMGENNCNTUFORCEUPMETHOD}\item [MCPWM\_GEN\regindex{n}\_CNTUFORCE\_UPMETHOD] 配置生成器\regindex{n} 的连续软件强制事件更新方式。\\
所有 bit 为 0 时：立即更新\\
bit0 为 1：发生 TEZ 事件时更新\\
bit1 为 1：发生 TEP 事件时更新\\
bit2 为 1：发生 TEA 事件时更新\\
bit3 为 1：发生 TEB 事件时更新\\
bit4：发生同步事件时更新\\
bit5 为 1：关闭更新\\
本文档中，TEA/B 指定时器值为寄存器 A/B 的值时生成的事件。\\
(R/W)
\label{fielddesc:MCPWMGENNACNTUFORCEMODE}\item [MCPWM\_GEN\regindex{n}\_A\_CNTUFORCE\_MODE] 配置 PWM\regindex{n} A 的连续软件强制模式。\\
0：关闭\\
1：低电平\\
2：高电平\\
3：关闭\\
(R/W)
\label{fielddesc:MCPWMGENNBCNTUFORCEMODE}\item [MCPWM\_GEN\regindex{n}\_B\_CNTUFORCE\_MODE] 配置 PWM\regindex{n} B 的连续软件强制模式。\\
0：关闭\\
1：低电平\\
2：高电平\\
3：关闭\\
(R/W)
\label{fielddesc:MCPWMGENNANCIFORCE}\item [MCPWM\_GEN\regindex{n}\_A\_NCIFORCE] 配置是否触发 PWM\regindex{n} A 上的非连续即时软件强制事件。\\
0：无效\\
1：触发 PWM\regindex{n} A 上的非连续即时软件强制事件\\
(R/W)
\item [见下页...]
\end{reglist}\end{regdesc}
\end{register}
\addtocounter{Regfloat}{-1}
\begin{register}{H}{MCPWM\_GEN\regindex{n}\_FORCE\_REG(\regindex{n}: 0-2)}
{0x004C+0x38*\regindex{n}}
\begin{regdesc}\begin{reglist}
\item [...接上页]
\label{fielddesc:MCPWMGENNANCIFORCEMODE}\item [MCPWM\_GEN\regindex{n}\_A\_NCIFORCE\_MODE] 配置用于 PWM\regindex{n} A 的非连续即时软件强制模式。\\
0：关闭\\
1：低电平\\
2：高电平\\
3：关闭\\
(R/W)
\label{fielddesc:MCPWMGENNBNCIFORCE}\item [MCPWM\_GEN\regindex{n}\_B\_NCIFORCE] 配置是否触发 PWM\regindex{n} B 上的非连续即时软件强制事件。\\
0：无效\\
1：触发 PWM\regindex{n} B 上的非连续即时软件强制事件\\
(R/W)
\label{fielddesc:MCPWMGENNBNCIFORCEMODE}\item [MCPWM\_GEN\regindex{n}\_B\_NCIFORCE\_MODE] 配置用于 PWM\regindex{n} B 的非连续即时软件强制模式。\\
0：关闭\\
1：低电平\\
2：高电平\\
3：关闭\\
(R/W)
\end{reglist}\end{regdesc}
\end{register}

\begin{register}{H}{MCPWM\_GEN\regindex{n}\_A\_REG
(\regindex{n}: 0-2)}
{0x0050+0x38*\regindex{n}}\label{regdesc:MCPWMGENNAREG}
\regfield{(reserved)}{8}{24}{00000000}%
\regfield{MCPWM\_GEN\regindex{n}\_A\_DT1}{2}{22}{{0}}%
\regfield{MCPWM\_GEN\regindex{n}\_A\_DT0}{2}{20}{{0}}%
\regfield{MCPWM\_GEN\regindex{n}\_A\_DTEB}{2}{18}{{0}}%
\regfield{MCPWM\_GEN\regindex{n}\_A\_DTEA}{2}{16}{{0}}%
\regfield{MCPWM\_GEN\regindex{n}\_A\_DTEP}{2}{14}{{0}}%
\regfield{MCPWM\_GEN\regindex{n}\_A\_DTEZ}{2}{12}{{0}}%
\regfield{MCPWM\_GEN\regindex{n}\_A\_UT1}{2}{10}{{0}}%
\regfield{MCPWM\_GEN\regindex{n}\_A\_UT0}{2}{8}{{0}}%
\regfield{MCPWM\_GEN\regindex{n}\_A\_UTEB}{2}{6}{{0}}%
\regfield{MCPWM\_GEN\regindex{n}\_A\_UTEA}{2}{4}{{0}}%
\regfield{MCPWM\_GEN\regindex{n}\_A\_UTEP}{2}{2}{{0}}%
\regfield{MCPWM\_GEN\regindex{n}\_A\_UTEZ}{2}{0}{{0}}%
\reglabel{Reset}\regnewline%
\begin{regdesc}\begin{reglist}
\label{fielddesc:MCPWMGENNAUTEZ}\item [MCPWM\_GEN\regindex{n}\_A\_UTEZ] 配置定时器递增时，TEZ 事件在 PWM\regindex{n}A 上触发的操作。\\0：波形无改变\\1：拉低\\2：拉高\\3：取反\\(R/W)
\label{fielddesc:MCPWMGENNAUTEP}\item [MCPWM\_GEN\regindex{n}\_A\_UTEP] 配置定时器递增时，TEP 事件在 PWM\regindex{n}A 上触发的操作。\\0：波形无改变\\1：拉低\\2：拉高\\3：取反\\(R/W)
\label{fielddesc:MCPWMGENNAUTEA}\item [MCPWM\_GEN\regindex{n}\_A\_UTEA] 配置定时器递增时，TEA 事件在 PWM\regindex{n}A 上触发的操作。\\0：波形无改变\\1：拉低\\2：拉高\\3：取反\\(R/W)
\label{fielddesc:MCPWMGENNAUTEB}\item [MCPWM\_GEN\regindex{n}\_A\_UTEB] 配置定时器递增时，TEB 事件在 PWM\regindex{n}A 上触发的操作。\\0：波形无改变\\1：拉低\\2：拉高\\3：取反\\(R/W)
\item [见下页...]
\end{reglist}\end{regdesc}
\end{register}
\addtocounter{Regfloat}{-1}
\begin{register}{H}{MCPWM\_GEN\regindex{n}\_A\_REG(\regindex{n}: 0-2)}{0x0050+0x38*\regindex{n}}
\begin{regdesc}\begin{reglist}
\item [...接上页]
\label{fielddesc:MCPWMGENNAUT0}\item [MCPWM\_GEN\regindex{n}\_A\_UT0] 配置定时器递增时，event\_t0 在 PWM\regindex{n}A 上触发的操作。\\0：波形无改变\\1：拉低\\2：拉高\\3：取反\\(R/W)
\label{fielddesc:MCPWMGENNAUT1}\item [MCPWM\_GEN\regindex{n}\_A\_UT1] 配置定时器递增时，event\_t1 在 PWM\regindex{n}A 上触发的操作。\\0：波形无改变\\1：拉低\\2：拉高\\3：取反\\(R/W)
\label{fielddesc:MCPWMGENNADTEZ}\item [MCPWM\_GEN\regindex{n}\_A\_DTEZ] 配置定时器递减时，TEZ 事件在 PWM\regindex{n}A 上触发的操作。\\0：波形无改变\\1：拉低\\2：拉高\\3：取反\\(R/W)
\label{fielddesc:MCPWMGENNADTEP}\item [MCPWM\_GEN\regindex{n}\_A\_DTEP] 配置定时器递减时，TEP 事件在 PWM\regindex{n}A 上触发的操作。\\0：波形无改变\\1：拉低\\2：拉高\\3：取反\\(R/W)
\label{fielddesc:MCPWMGENNADTEA}\item [MCPWM\_GEN\regindex{n}\_A\_DTEA] 配置定时器递减时，TEA 事件在 PWM\regindex{n}A 上触发的操作。\\0：波形无改变\\1：拉低\\2：拉高\\3：取反\\(R/W)
\item [见下页...]
\end{reglist}\end{regdesc}
\end{register}
\addtocounter{Regfloat}{-1}
\begin{register}{H}{MCPWM\_GEN\regindex{n}\_A\_REG(\regindex{n}: 0-2)}{0x0050+0x38*\regindex{n}}
\begin{regdesc}\begin{reglist}
\item [...接上页]
\label{fielddesc:MCPWMGENNADTEB}\item [MCPWM\_GEN\regindex{n}\_A\_DTEB] 配置定时器递减时，TEB 事件在 PWM\regindex{n}A 上触发的操作。\\0：波形无改变\\1：拉低\\2：拉高\\3：取反\\(R/W)
\label{fielddesc:MCPWMGENNADT0}\item [MCPWM\_GEN\regindex{n}\_A\_DT0] 配置定时器递减时，event\_t0 在 PWM\regindex{n}A 上触发的操作。\\0：波形无改变\\1：拉低\\2：拉高\\3：取反\\(R/W)
\label{fielddesc:MCPWMGENNADT1}\item [MCPWM\_GEN\regindex{n}\_A\_DT1] 配置定时器递减时，event\_t1 在 PWM\regindex{n}A 上触发的操作。\\0：波形无改变\\1：拉低\\2：拉高\\3：取反\\(R/W)
\end{reglist}\end{regdesc}
\end{register}


\begin{register}{H}{MCPWM\_GEN\regindex{n}\_B\_REG
(\regindex{n}: 0-2)}
{0x0054+0x38*\regindex{n}}\label{regdesc:MCPWMGENNBREG}
\regfield{(reserved)}{8}{24}{00000000}%
\regfield{MCPWM\_GEN\regindex{n}\_B\_DT1}{2}{22}{{0}}%
\regfield{MCPWM\_GEN\regindex{n}\_B\_DT0}{2}{20}{{0}}%
\regfield{MCPWM\_GEN\regindex{n}\_B\_DTEB}{2}{18}{{0}}%
\regfield{MCPWM\_GEN\regindex{n}\_B\_DTEA}{2}{16}{{0}}%
\regfield{MCPWM\_GEN\regindex{n}\_B\_DTEP}{2}{14}{{0}}%
\regfield{MCPWM\_GEN\regindex{n}\_B\_DTEZ}{2}{12}{{0}}%
\regfield{MCPWM\_GEN\regindex{n}\_B\_UT1}{2}{10}{{0}}%
\regfield{MCPWM\_GEN\regindex{n}\_B\_UT0}{2}{8}{{0}}%
\regfield{MCPWM\_GEN\regindex{n}\_B\_UTEB}{2}{6}{{0}}%
\regfield{MCPWM\_GEN\regindex{n}\_B\_UTEA}{2}{4}{{0}}%
\regfield{MCPWM\_GEN\regindex{n}\_B\_UTEP}{2}{2}{{0}}%
\regfield{MCPWM\_GEN\regindex{n}\_B\_UTEZ}{2}{0}{{0}}%
\reglabel{Reset}\regnewline%
\begin{regdesc}\begin{reglist}
\label{fielddesc:MCPWMGENNBUTEZ}\item [MCPWM\_GEN\regindex{n}\_B\_UTEZ] 定时器递增时，TEZ 在 PWM\regindex{n}B 上触发的操作。\\0：波形无改变\\1：拉低\\2：拉高\\3：取反\\ (R/W)
\label{fielddesc:MCPWMGENNBUTEP}\item [MCPWM\_GEN\regindex{n}\_B\_UTEP] 定时器递增时，TEP 在 PWM\regindex{n}B 上触发的操作。\\0：波形无改变\\1：拉低\\2：拉高\\3：取反\\ (R/W)
\label{fielddesc:MCPWMGENNBUTEA}\item [MCPWM\_GEN\regindex{n}\_B\_UTEA] 定时器递增时，TEA 在 PWM\regindex{n}B 上触发的操作。\\0：波形无改变\\1：拉低\\2：拉高\\3：取反\\ (R/W)
\label{fielddesc:MCPWMGENNBUTEB}\item [MCPWM\_GEN\regindex{n}\_B\_UTEB] 定时器递增时，TEB 在 PWM\regindex{n}B 上触发的操作。\\0：波形无改变\\1：拉低\\2：拉高\\3：取反\\ (R/W)
\item [见下页...]
\end{reglist}\end{regdesc}
\end{register}
\addtocounter{Regfloat}{-1}
\begin{register}{H}{MCPWM\_GEN\regindex{n}\_B\_REG(\regindex{n}: 0-2)}{0x0054+0x38*\regindex{n}}
\begin{regdesc}\begin{reglist}
\item [...接上页]
\label{fielddesc:MCPWMGENNBUT0}\item [MCPWM\_GEN\regindex{n}\_B\_UT0] 定时器递增时，event\_t0 在 PWM\regindex{n}B 上触发的操作。\\0：波形无改变\\1：拉低\\2：拉高\\3：取反\\ (R/W)
\label{fielddesc:MCPWMGENNBUT1}\item [MCPWM\_GEN\regindex{n}\_B\_UT1] 定时器递增时，event\_t1 在 PWM\regindex{n}B 上触发的操作。\\0：波形无改变\\1：拉低\\2：拉高\\3：取反\\ (R/W)
\label{fielddesc:MCPWMGENNBDTEZ}\item [MCPWM\_GEN\regindex{n}\_B\_DTEZ] 定时器递减时，TEZ 事件在 PWM\regindex{n}B 上触发的操作。\\0：波形无改变\\1：拉低\\2：拉高\\3：取反\\ (R/W)
\label{fielddesc:MCPWMGENNBDTEP}\item [MCPWM\_GEN\regindex{n}\_B\_DTEP] 定时器递减时，TEP 事件在 PWM\regindex{n}B 上触发的操作。\\0：波形无改变\\1：拉低\\2：拉高\\3：取反\\ (R/W)
\label{fielddesc:MCPWMGENNBDTEA}\item [MCPWM\_GEN\regindex{n}\_B\_DTEA] 定时器递减时，TEA 事件在 PWM\regindex{n}B 上触发的操作。\\0：波形无改变\\1：拉低\\2：拉高\\3：取反\\ (R/W)
\item [见下页...]
\end{reglist}\end{regdesc}
\end{register}
\addtocounter{Regfloat}{-1}
\begin{register}{H}{MCPWM\_GEN\regindex{n}\_B\_REG(\regindex{n}: 0-2)}{0x0054+0x38*\regindex{n}}
\begin{regdesc}\begin{reglist}
\item [...接上页]
\label{fielddesc:MCPWMGENNBDTEB}\item [MCPWM\_GEN\regindex{n}\_B\_DTEB] 定时器递减时，TEB 事件在 PWM\regindex{n}B 上触发的操作。\\0：波形无改变\\1：拉低\\2：拉高\\3：取反\\ (R/W)
\label{fielddesc:MCPWMGENNBDT0}\item [MCPWM\_GEN\regindex{n}\_B\_DT0] 定时器递减时，event\_t0 事件在 PWM\regindex{n}B 上触发的操作。\\0：波形无改变\\1：拉低\\2：拉高\\3：取反\\ (R/W)
\label{fielddesc:MCPWMGENNBDT1}\item [MCPWM\_GEN\regindex{n}\_B\_DT1] 定时器递减时，event\_t1 在 PWM\regindex{n}B 上触发的操作。\\0：波形无改变\\1：拉低\\2：拉高\\3：取反\\ (R/W)
\end{reglist}\end{regdesc}
\end{register}


\begin{register}{H}{MCPWM\_DT\regindex{n}\_CFG\_REG(\regindex{n}: 0-2)}
{0x0058+0x38*\regindex{n}}\label{regdesc:MCPWMDTNCFGREG}
\regfield{(reserved)}{14}{18}{00000000000000}%
\regfield{MCPWM\_DB\regindex{n}\_CLK\_SEL}{1}{17}{{0}}%
\regfield{MCPWM\_DB\regindex{n}\_B\_OUTBYPASS}{1}{16}{{1}}%
\regfield{MCPWM\_DB\regindex{n}\_A\_OUTBYPASS}{1}{15}{{1}}%
\regfield{MCPWM\_DB\regindex{n}\_FED\_OUTINVERT}{1}{14}{{0}}%
\regfield{MCPWM\_DB\regindex{n}\_RED\_OUTINVERT}{1}{13}{{0}}%
\regfield{MCPWM\_DB\regindex{n}\_FED\_INSEL}{1}{12}{{0}}%
\regfield{MCPWM\_DB\regindex{n}\_RED\_INSEL}{1}{11}{{0}}%
\regfield{MCPWM\_DB\regindex{n}\_B\_OUTSWAP}{1}{10}{{0}}%
\regfield{MCPWM\_DB\regindex{n}\_A\_OUTSWAP}{1}{9}{{0}}%
\regfield{MCPWM\_DB\regindex{n}\_DEB\_MODE}{1}{8}{{0}}%
\regfield{MCPWM\_DB\regindex{n}\_RED\_UPMETHOD}{4}{4}{{0}}%
\regfield{MCPWM\_DB\regindex{n}\_FED\_UPMETHOD}{4}{0}{{0}}%
\reglabel{Reset}\regnewline%
\begin{regdesc}\begin{reglist}
\label{fielddesc:MCPWMDBNFEDUPMETHOD}\item [MCPWM\_DB\regindex{n}\_FED\_UPMETHOD] 配置下降沿延迟有效寄存器的更新方式。\\
所有位值为 0：立即更新\\
bit0 为 1：发生 TEZ 事件时更新\\
bit1 为 1：发生 TEP 事件时更新\\
bit2 为 1：发生同步时间时更新\\
bit3 为 1：关闭更新\\
(R/W)
\label{fielddesc:MCPWMDBNREDUPMETHOD}\item [MCPWM\_DB\regindex{n}\_RED\_UPMETHOD] 配置上升沿延迟有效寄存器的更新方式。\\
所有位值为 0：立即更新\\
bit0 为 1：发生 TEZ 事件时更新\\
bit1 为 1：发生 TEP 事件时更新\\
bit2 为 1：发生同步时间时更新\\
bit3 为 1：关闭更新\\
(R/W)
\label{fielddesc:MCPWMDBNDEBMODE}\item [MCPWM\_DB\regindex{n}\_DEB\_MODE] 配置表 \ref{tab:mcpwm-s-dt}中的 S8 开关。详细配置信息请参考表 \ref{tab:mcpwm-mod-dt}。\\
0：FED/RED 分别在不同路径上生效\\
1：FED/RED 在 B 路径上生效，A 输出为旁路或 dulpB 模式 \\
(R/W)
\label{fielddesc:MCPWMDBNAOUTSWAP}\item [MCPWM\_DB\regindex{n}\_A\_OUTSWAP] 配置表 \ref{tab:mcpwm-s-dt}中的 S6 开关。详细配置信息请参考表 \ref{tab:mcpwm-mod-dt}。 (R/W)
\label{fielddesc:MCPWMDBNBOUTSWAP}\item [MCPWM\_DB\regindex{n}\_B\_OUTSWAP] 配置表 \ref{tab:mcpwm-s-dt}中的 S7 开关。详细配置信息请参考表 \ref{tab:mcpwm-mod-dt}。 (R/W)
\label{fielddesc:MCPWMDBNREDINSEL}\item [MCPWM\_DB\regindex{n}\_RED\_INSEL] 配置表 \ref{tab:mcpwm-s-dt}中的 S4 开关。详细配置信息请参考表 \ref{tab:mcpwm-mod-dt}。 (R/W)
\label{fielddesc:MCPWMDBNFEDINSEL}\item [MCPWM\_DB\regindex{n}\_FED\_INSEL] 配置表 \ref{tab:mcpwm-s-dt}中的 S5 开关。详细配置信息请参考表 \ref{tab:mcpwm-mod-dt}。 (R/W)
\label{fielddesc:MCPWMDBNREDOUTINVERT}\item [MCPWM\_DB\regindex{n}\_RED\_OUTINVERT] 配置表 \ref{tab:mcpwm-s-dt}中的 S2 开关。详细配置信息请参考表 \ref{tab:mcpwm-mod-dt}。 (R/W)
\label{fielddesc:MCPWMDBNFEDOUTINVERT}\item [MCPWM\_DB\regindex{n}\_FED\_OUTINVERT] 配置表 \ref{tab:mcpwm-s-dt}中的 S3 开关。详细配置信息请参考表 \ref{tab:mcpwm-mod-dt}。 (R/W)
\label{fielddesc:MCPWMDBNAOUTBYPASS}\item [MCPWM\_DB\regindex{n}\_A\_OUTBYPASS] 配置表 \ref{tab:mcpwm-s-dt}中的 S1 开关。详细配置信息请参考表 \ref{tab:mcpwm-mod-dt}。 (R/W)
\label{fielddesc:MCPWMDBNBOUTBYPASS}\item [MCPWM\_DB\regindex{n}\_B\_OUTBYPASS] 配置表 \ref{tab:mcpwm-s-dt}中的 S0 开关。详细配置信息请参考表 \ref{tab:mcpwm-mod-dt}。 (R/W)
\item [见下页...]
\end{reglist}\end{regdesc}
\end{register}
\addtocounter{Regfloat}{-1}
\begin{register}{H}{MCPWM\_DT\regindex{n}\_CFG\_REG(\regindex{n}: 0-2)}{0x0058+0x38*\regindex{n}}
\begin{regdesc}\begin{reglist}
\item [...接上页]
\label{fielddesc:MCPWMDBNCLKSEL}\item [MCPWM\_DB\regindex{n}\_CLK\_SEL] 选择死区时间生成器 0 的时钟。\\0：PWM\_CLK\\1：PT\_CLK\\ (R/W)
\end{reglist}\end{regdesc}
\end{register}


\begin{register}{H}{MCPWM\_DT\regindex{n}\_FED\_CFG\_REG(\regindex{n}: 0-2)}
{0x005C+0x38*\regindex{n}}\label{regdesc:MCPWMDTNFEDCFGREG}
\regfield{(reserved)}{16}{16}{0000000000000000}%
\regfield{MCPWM\_DB\regindex{n}\_FED}{16}{0}{{0}}%
\reglabel{Reset}\regnewline%
\begin{regdesc}\begin{reglist}
\label{fielddesc:MCPWMDBNFED}\item [MCPWM\_DB\regindex{n}\_FED] 配置 FED 的影子寄存器。 (R/W)
\end{reglist}\end{regdesc}
\end{register}


\begin{register}{H}{MCPWM\_DT\regindex{n}\_RED\_CFG\_REG(\regindex{n}: 0-2)}
{0x0060+0x38*\regindex{n}}\label{regdesc:MCPWMDTNREDCFGREG}
\regfield{(reserved)}{16}{16}{0000000000000000}%
\regfield{MCPWM\_DB\regindex{n}\_RED}{16}{0}{{0}}%
\reglabel{Reset}\regnewline%
\begin{regdesc}\begin{reglist}
\label{fielddesc:MCPWMDBNRED}\item [MCPWM\_DB\regindex{n}\_RED] 配置 RED 的影子寄存器。(R/W)
\end{reglist}\end{regdesc}
\end{register}


\begin{register}{H}{MCPWM\_CARRIER\regindex{n}\_CFG\_REG(\regindex{n}: 0-2)}
{0x0064+0x38*\regindex{n}}\label{regdesc:MCPWMCARRIERNCFGREG}
\regfield{(reserved)}{18}{14}{000000000000000000}%
\regfield{MCPWM\_CHOPPER\regindex{n}\_IN\_INVERT}{1}{13}{{0}}%
\regfield{MCPWM\_CHOPPER\regindex{n}\_OUT\_INVERT}{1}{12}{{0}}%
\regfield{MCPWM\_CHOPPER\regindex{n}\_OSHTWTH}{4}{8}{{0}}%
\regfield{MCPWM\_CHOPPER\regindex{n}\_DUTY}{3}{5}{{0}}%
\regfield{MCPWM\_CHOPPER\regindex{n}\_PRESCALE}{4}{1}{{0}}%
\regfield{MCPWM\_CHOPPER\regindex{n}\_EN}{1}{0}{{0}}%
\reglabel{Reset}\regnewline%
\begin{regdesc}\begin{reglist}
\label{fielddesc:MCPWMCHOPPERNEN}\item [MCPWM\_CHOPPER\regindex{n}\_EN] 配置是否使能载波\regindex{n}。\\
0：绕过\\
1：使能\\
(R/W)
\label{fielddesc:MCPWMCHOPPERNPRESCALE}\item [MCPWM\_CHOPPER\regindex{n}\_PRESCALE] 配置 PWM 载波\regindex{n} 时钟 (PC\_CLK) 的预分频值。PC\_CLK 周期 = PWM\_CLK 周期 * (PWM\_CARRIER\regindex{n}\_PRESCALE + 1)。 (R/W)
\label{fielddesc:MCPWMCHOPPERNDUTY}\item [MCPWM\_CHOPPER\regindex{n}\_DUTY] 配置载波占空比。占空比 = PWM\_CARRIER\regindex{n}\_DUTY/8。 (R/W)
\label{fielddesc:MCPWMCHOPPERNOSHTWTH}\item [MCPWM\_CHOPPER\regindex{n}\_OSHTWTH] 配置载波第一个脉冲的宽度，单位为载波周期。 (R/W)
\label{fielddesc:MCPWMCHOPPERNOUTINVERT}\item [MCPWM\_CHOPPER\regindex{n}\_OUT\_INVERT] 配置是否将此模块的 PWM\regindex{n}A 和 PWM\regindex{n}B 输出反相。\\
0：正常输出\\
1：反相\\
(R/W)
\label{fielddesc:MCPWMCHOPPERNININVERT}\item [MCPWM\_CHOPPER\regindex{n}\_IN\_INVERT] 配置是否将此模块的 PWM\regindex{n}A 和 PWM\regindex{n}B 输入反相。\\
0：正常输入\\
1：反相\\
(R/W)
\end{reglist}\end{regdesc}
\end{register}


\begin{register}{H}{MCPWM\_FH\regindex{n}\_CFG0\_REG(\regindex{n}: 0-2)}
{0x0068+0x38*\regindex{n}}\label{regdesc:MCPWMFHNCFG0REG}
\regfield{(reserved)}{8}{24}{00000000}%
\regfield{MCPWM\_TZ\regindex{n}\_B\_OST\_U}{2}{22}{{0}}%
\regfield{MCPWM\_TZ\regindex{n}\_B\_OST\_D}{2}{20}{{0}}%
\regfield{MCPWM\_TZ\regindex{n}\_B\_CBC\_U}{2}{18}{{0}}%
\regfield{MCPWM\_TZ\regindex{n}\_B\_CBC\_D}{2}{16}{{0}}%
\regfield{MCPWM\_TZ\regindex{n}\_A\_OST\_U}{2}{14}{{0}}%
\regfield{MCPWM\_TZ\regindex{n}\_A\_OST\_D}{2}{12}{{0}}%
\regfield{MCPWM\_TZ\regindex{n}\_A\_CBC\_U}{2}{10}{{0}}%
\regfield{MCPWM\_TZ\regindex{n}\_A\_CBC\_D}{2}{8}{{0}}%
\regfield{MCPWM\_TZ\regindex{n}\_F\regindex{n}\_OST}{1}{7}{{0}}%
\regfield{MCPWM\_TZ\regindex{n}\_F1\_OST}{1}{6}{{0}}%
\regfield{MCPWM\_TZ\regindex{n}\_F2\_OST}{1}{5}{{0}}%
\regfield{MCPWM\_TZ\regindex{n}\_SW\_OST}{1}{4}{{0}}%
\regfield{MCPWM\_TZ\regindex{n}\_F0\_CBC}{1}{3}{{0}}%
\regfield{MCPWM\_TZ\regindex{n}\_F1\_CBC}{1}{2}{{0}}%
\regfield{MCPWM\_TZ\regindex{n}\_F2\_CBC}{1}{1}{{0}}%
\regfield{MCPWM\_TZ\regindex{n}\_SW\_CBC}{1}{0}{{0}}%
\reglabel{Reset}\regnewline%
\begin{regdesc}\begin{reglist}
\label{fielddesc:MCPWMTZNSWCBC}\item [MCPWM\_TZ\regindex{n}\_SW\_CBC] 配置是否使能软件强制逐周期模式操作。\\0：关闭\\1：使能\\ (R/W)
\label{fielddesc:MCPWMTZNF2CBC}\item [MCPWM\_TZ\regindex{n}\_F2\_CBC] 配置 fault\_event2 是否触发逐周期模式操作。\\0：关闭\\1：使能\\ (R/W)
\label{fielddesc:MCPWMTZNF1CBC}\item [MCPWM\_TZ\regindex{n}\_F1\_CBC] 配置 fault\_event1 是否触发逐周期模式操作。\\0：关闭\\1：使能\\ (R/W)
\label{fielddesc:MCPWMTZNF0CBC}\item [MCPWM\_TZ\regindex{n}\_F0\_CBC] 配置 fault\_event0 是否触发逐周期模式操作。\\0：关闭\\1：使能\\ (R/W)
\label{fielddesc:MCPWMTZNSWOST}\item [MCPWM\_TZ\regindex{n}\_SW\_OST] 配置是否使能软件强制一次性模式操作。\\0：关闭\\1：使能\\ (R/W)
\label{fielddesc:MCPWMTZNF2OST}\item [MCPWM\_TZ\regindex{n}\_F2\_OST] 配置 fault\_event2 是否触发一次性模式操作。\\0：关闭\\1：使能\\ (R/W)
\item [见下页...]
\end{reglist}\end{regdesc}
\end{register}
\addtocounter{Regfloat}{-1}
\begin{register}{H}{MCPWM\_FH\regindex{n}\_CFG0\_REG(\regindex{n}: 0-2)}{0x0068+0x38*\regindex{n}}
\begin{regdesc}\begin{reglist}
\item [...接上页]
\label{fielddesc:MCPWMTZNF1OST}\item [MCPWM\_TZ\regindex{n}\_F1\_OST] 配置 fault\_event1 是否触发一次性模式操作。\\0：关闭\\1：使能\\ (R/W)
\label{fielddesc:MCPWMTZNFNOST}\item [MCPWM\_TZ\regindex{n}\_F\regindex{n}\_OST] 配置 fault\_event0 是否触发一次性模式操作。\\0：关闭\\1：使能\\ (R/W)
\label{fielddesc:MCPWMTZNACBCD}\item [MCPWM\_TZ\regindex{n}\_A\_CBC\_D] 配置定时器递减计数并且发生故障事件时，PWM\regindex{n}A 上的逐周期模式操作。\\0：无\\1：强制拉低\\2：强制拉高\\3：取反\\ (R/W)
\label{fielddesc:MCPWMTZNACBCU}\item [MCPWM\_TZ\regindex{n}\_A\_CBC\_U] 配置定时器递增计数并且发生故障事件时，PWM\regindex{n}A 上的逐周期模式操作。\\0：无\\1：强制拉低\\2：强制拉高\\3：取反\\ (R/W)
\label{fielddesc:MCPWMTZNAOSTD}\item [MCPWM\_TZ\regindex{n}\_A\_OST\_D] 配置定时器递减计数并且发生故障事件时，PWM\regindex{n}A 上的一次性模式操作。\\0：无\\1：强制拉低\\2：强制拉高\\3：取反\\ (R/W)
\item [见下页...]
\end{reglist}\end{regdesc}
\end{register}
\addtocounter{Regfloat}{-1}
\begin{register}{H}{MCPWM\_FH\regindex{n}\_CFG0\_REG(\regindex{n}: 0-2)}{0x0068+0x38*\regindex{n}}
\begin{regdesc}\begin{reglist}
\item [...接上页]
\label{fielddesc:MCPWMTZNAOSTU}\item [MCPWM\_TZ\regindex{n}\_A\_OST\_U] 配置定时器递增计数并且发生故障事件时，PWM\regindex{n}A 上的一次性模式操作。\\0：无\\1：强制拉低\\2：强制拉高\\3：取反\\ (R/W)
\label{fielddesc:MCPWMTZNBCBCD}\item [MCPWM\_TZ\regindex{n}\_B\_CBC\_D] 配置定时器递减计数并且发生故障事件时，PWM\regindex{n}B 上的逐周期模式操作。\\0：无\\1：强制拉低\\2：强制拉高\\3：取反\\ (R/W)
\label{fielddesc:MCPWMTZNBCBCU}\item [MCPWM\_TZ\regindex{n}\_B\_CBC\_U] 配置定时器递增计数并且发生故障事件时，PWM\regindex{n}B 上的逐周期模式操作。\\0：无\\1：强制拉低\\2：强制拉高\\3：取反\\ (R/W)
\label{fielddesc:MCPWMTZNBOSTD}\item [MCPWM\_TZ\regindex{n}\_B\_OST\_D] 配置定时器递减计数并且发生故障事件时，PWM\regindex{n}B 上的一次性模式操作。\\0：无\\1：强制拉低\\2：强制拉高\\3：取反\\ (R/W)
\label{fielddesc:MCPWMTZNBOSTU}\item [MCPWM\_TZ\regindex{n}\_B\_OST\_U] 配置定时器递增计数并且发生故障事件时，PWM\regindex{n}B 上的一次性模式操作。\\0：无\\1：强制拉低\\2：强制拉高\\3：取反\\ (R/W)
\end{reglist}\end{regdesc}
\end{register}


\begin{register}{H}{MCPWM\_FH\regindex{n}\_CFG1\_REG(\regindex{n}: 0-2)}
{0x006C+0x38*\regindex{n}}\label{regdesc:MCPWMFHNCFG1REG}
\regfield{(reserved)}{27}{5}{000000000000000000000000000}%
\regfield{MCPWM\_TZ\regindex{n}\_FORCE\_OST}{1}{4}{{0}}%
\regfield{MCPWM\_TZ\regindex{n}\_FORCE\_CBC}{1}{3}{{0}}%
\regfield{MCPWM\_TZ\regindex{n}\_CBCPULSE}{2}{1}{{0}}%
\regfield{MCPWM\_TZ\regindex{n}\_CLR\_OST}{1}{0}{{0}}%
\reglabel{Reset}\regnewline%
\begin{regdesc}\begin{reglist}
\label{fielddesc:MCPWMTZNCLROST}\item [MCPWM\_TZ\regindex{n}\_CLR\_OST] 配置是否通过软件清除持续一次性模式操作。\\
0：无效\\
1：通过软件清除持续一次性模式操作\\
(R/W)
\label{fielddesc:MCPWMTZNCBCPULSE}\item [MCPWM\_TZ\regindex{n}\_CBCPULSE] 配置逐周期模式操作更新的时间点。\\
bit0 为 1：发生 TEZ 事件时\\
bit1 为 1：发生 TEP 事件时\\
(R/W)
\label{fielddesc:MCPWMTZNFORCECBC}\item [MCPWM\_TZ\regindex{n}\_FORCE\_CBC] 配置是否触发逐周期模式的操作。\\
0：无效\\
1：触发\\
(R/W)
\label{fielddesc:MCPWMTZNFORCEOST}\item [MCPWM\_TZ\regindex{n}\_FORCE\_OST] 配置是否触发一次性模式的操作。\\
0：无效\\
1：触发\\
(R/W)
\end{reglist}\end{regdesc}
\end{register}


\begin{register}{H}{MCPWM\_FH\regindex{n}\_STATUS\_REG(\regindex{n}: 0-2)}{0x0070+0x38*\regindex{n}}\label{regdesc:MCPWMFHNSTATUSREG}
\regfield{(reserved)}{30}{2}{000000000000000000000000000000}%
\regfield{MCPWM\_TZ\regindex{n}\_OST\_ON}{1}{1}{{0}}%
\regfield{MCPWM\_TZ\regindex{n}\_CBC\_ON}{1}{0}{{0}}%
\reglabel{Reset}\regnewline%
\begin{regdesc}\begin{reglist}
\label{fielddesc:MCPWMTZNCBCON}\item [MCPWM\_TZ\regindex{n}\_CBC\_ON] 表示是否正在进行逐周期模式的操作。\\
0：不存在正在进行的逐周期模式操作\\
1：逐周期模式的操作正在进行\\
(RO)
\label{fielddesc:MCPWMTZNOSTON}\item [MCPWM\_TZ\regindex{n}\_OST\_ON] 表示是否正在进行一次性模式的操作。\\
0：不存在正在进行的一次性模式操作\\
1：一次性模式的操作正在进行\\
(RO)
\end{reglist}\end{regdesc}
\end{register}


\begin{register}{H}{MCPWM\_FAULT\_DETECT\_REG}{0x00E4}\label{regdesc:MCPWMFAULTDETECTREG}
\regfield{(reserved)}{23}{9}{00000000000000000000000}%
\regfield{MCPWM\_EVENT\_F2}{1}{8}{{0}}%
\regfield{MCPWM\_EVENT\_F1}{1}{7}{{0}}%
\regfield{MCPWM\_EVENT\_F0}{1}{6}{{0}}%
\regfield{MCPWM\_F2\_POLE}{1}{5}{{0}}%
\regfield{MCPWM\_F1\_POLE}{1}{4}{{0}}%
\regfield{MCPWM\_F0\_POLE}{1}{3}{{0}}%
\regfield{MCPWM\_F2\_EN}{1}{2}{{0}}%
\regfield{MCPWM\_F1\_EN}{1}{1}{{0}}%
\regfield{MCPWM\_F0\_EN}{1}{0}{{0}}%
\reglabel{Reset}\regnewline%
\begin{regdesc}\begin{reglist}
\label{fielddesc:MCPWMFNEN}\item [MCPWM\_F\regindex{n}\_EN] 配置是否使能生成 fault\_event\regindex{n}。\\
0：无效\\
1：使能\\
(R/W)
\label{fielddesc:MCPWMFNPOLE}\item [MCPWM\_F\regindex{n}\_POLE] 配置来自 GPIO 矩阵的 FAULT\regindex{n} 信号源触发 fault\_event\regindex{n} 时极性。\\
0：低电平触发\\
1：高电平触发\\
(R/W)
\label{fielddesc:MCPWMEVENTFN}\item [MCPWM\_EVENT\_F\regindex{n}] 表示 fault\_event\regindex{n} 事件的状态。此字段由硬件置 1 和清零。\\
0：fault\_event\regindex{n} 事件停止\\
1：fault\_event\regindex{n} 事件持续\\
(RO)
\end{reglist}\end{regdesc}
\end{register}


\begin{register}{H}{MCPWM\_CAP\_TIMER\_CFG\_REG}{0x00E8}\label{regdesc:MCPWMCAPTIMERCFGREG}
\regfield{(reserved)}{26}{6}{00000000000000000000000000}%
\regfield{MCPWM\_CAP\_SYNC\_SW}{1}{5}{{0}}%
\regfield{MCPWM\_CAP\_SYNCI\_SEL}{3}{2}{{0}}%
\regfield{MCPWM\_CAP\_SYNCI\_EN}{1}{1}{{0}}%
\regfield{MCPWM\_CAP\_TIMER\_EN}{1}{0}{{0}}%
\reglabel{Reset}\regnewline%
\begin{regdesc}\begin{reglist}
\label{fielddesc:MCPWMCAPTIMEREN}\item [MCPWM\_CAP\_TIMER\_EN] 配置是否使能捕获定时器在 APB\_CLK 下的递增。\\
0：无效\\
1：使能捕获定时器在 APB\_CLK 下的递增\\
(R/W)
\label{fielddesc:MCPWMCAPSYNCIEN}\item [MCPWM\_CAP\_SYNCI\_EN] 配置是否使能捕获定时器同步。\\
0：无效\\
1：使能捕获定时器同步\\
(R/W)
\label{fielddesc:MCPWMCAPSYNCISEL}\item [MCPWM\_CAP\_SYNCI\_SEL] 配置捕获模块的同步输入。\\
0：无\\
1：定时器 0 的同步输出\\
2：定时器 1 的同步输出\\
3：定时器 2 的同步输出\\
4：来自 GPIO 矩阵的 SYNC0\\
5：来自 GPIO 矩阵的 SYNC1\\
6：来自 GPIO 矩阵的 SYNC2\\
7：无\\
(R/W)
\label{fielddesc:MCPWMCAPSYNCSW}\item [MCPWM\_CAP\_SYNC\_SW] 当 \hyperref[fielddesc:MCPWMCAPSYNCIEN]{MCPWM\_CAP\_SYNCI\_EN} 置位时，配置是否同步捕获定时器，在捕获定时器中写入相位寄存器的值。\\
0：无效\\
1：同步捕获定时器\\
(WT)
\end{reglist}\end{regdesc}
\end{register}


\begin{register}{H}{MCPWM\_CAP\_TIMER\_PHASE\_REG}{0x00EC}\label{regdesc:MCPWMCAPTIMERPHASEREG}
\regfield{MCPWM\_CAP\_TIMER\_PHASE}{32}{0}{{0}}%
\reglabel{Reset}\regnewline%
\begin{regdesc}\begin{reglist}
\label{fielddesc:MCPWMCAPTIMERPHASE}\item [MCPWM\_CAP\_TIMER\_PHASE] 配置捕获定时器同步操作的相位值。(R/W)
\end{reglist}\end{regdesc}
\end{register}


\begin{register}{H}{MCPWM\_CAP\_CH\regindex{n}\_CFG\_REG (\regindex{n}: 0-2)}{0x00F0+0x4*\regindex{n}}\label{regdesc:MCPWMCAPCHNCFGREG}
\regfield{(reserved)}{19}{13}{0000000000000000000}%
\regfield{MCPWM\_CAP\regindex{n}\_SW}{1}{12}{{0}}%
\regfield{MCPWM\_CAP\regindex{n}\_IN\_INVERT}{1}{11}{{0}}%
\regfield{MCPWM\_CAP\regindex{n}\_PRESCALE}{8}{3}{{0}}%
\regfield{MCPWM\_CAP\regindex{n}\_MODE}{2}{1}{{0}}%
\regfield{MCPWM\_CAP\regindex{n}\_EN}{1}{0}{{0}}%
\reglabel{Reset}\regnewline%
\begin{regdesc}\begin{reglist}
\label{fielddesc:MCPWMCAPNEN}\item [MCPWM\_CAP\regindex{n}\_EN] 配置是否使能信道 \regindex{n} 上的捕获事件。\\
0：不使能\\
1：使能\\
(R/W)
\label{fielddesc:MCPWMCAPNMODE}\item [MCPWM\_CAP\regindex{n}\_MODE] 配置预分频后信道 \regindex{n} 上的捕获沿。\\
0：无\\
bit0 为 1：使能下降沿捕获。\\
bit1 为 1：使能上升沿捕获。 \\
(R/W)
\label{fielddesc:MCPWMCAPNPRESCALE}\item [MCPWM\_CAP\regindex{n}\_PRESCALE] 配置 CAP\regindex{n} 上升沿的预分频值。预分频值 = PWM\_CAP\regindex{n}\_PRESCALE + 1。 (R/W)
\label{fielddesc:MCPWMCAPNININVERT}\item [MCPWM\_CAP\regindex{n}\_IN\_INVERT] 配置是否在预分频前反相来自 GPIO 矩阵的 CAP\regindex{n}。\\
0：无效\\
1：在预分频前反相来自 GPIO 矩阵的 CAP\regindex{n}\\
(R/W)
\label{fielddesc:MCPWMCAPNSW}\item [MCPWM\_CAP\regindex{n}\_SW] 配置是否触发信道 0 上的软件强制捕获事件。\\
0：不触发\\
1：触发\\
(WT)
\end{reglist}\end{regdesc}
\end{register}


\begin{register}{H}{MCPWM\_CAP\_CH\regindex{n}\_REG (\regindex{n}: 0-2)}{0x00FC+0x4*\regindex{n}}\label{regdesc:MCPWMCAPCHNREG}
\regfield{MCPWM\_CAP\regindex{n}\_VALUE}{32}{0}{{0}}%
\reglabel{Reset}\regnewline%
\begin{regdesc}\begin{reglist}
\label{fielddesc:MCPWMCAPNVALUE}\item [MCPWM\_CAP\regindex{n}\_VALUE] 表示信道 \regindex{n} 上一次捕获的值。 (RO)
\end{reglist}\end{regdesc}
\end{register}


\begin{register}{H}{MCPWM\_CAP\_STATUS\_REG}{0x0108}\label{regdesc:MCPWMCAPSTATUSREG}
\regfield{(reserved)}{29}{3}{00000000000000000000000000000}%
\regfield{MCPWM\_CAP2\_EDGE}{1}{2}{{0}}%
\regfield{MCPWM\_CAP1\_EDGE}{1}{1}{{0}}%
\regfield{MCPWM\_CAP0\_EDGE}{1}{0}{{0}}%
\reglabel{Reset}\regnewline%
\begin{regdesc}\begin{reglist}
\label{fielddesc:MCPWMCAPNEDGE}\item [MCPWM\_CAP\regindex{n}\_EDGE] 信道\regindex{n} 上一次捕获触发事件的边沿。\\0：上升沿\\1：下降沿 \\(RO)
\end{reglist}\end{regdesc}
\end{register}


\begin{register}{H}{MCPWM\_UPDATE\_CFG\_REG}{0x010C}\label{regdesc:MCPWMUPDATECFGREG}
\regfield{(reserved)}{24}{8}{000000000000000000000000}%
\regfield{MCPWM\_OP2\_FORCE\_UP}{1}{7}{{0}}%
\regfield{MCPWM\_OP2\_UP\_EN}{1}{6}{{1}}%
\regfield{MCPWM\_OP1\_FORCE\_UP}{1}{5}{{0}}%
\regfield{MCPWM\_OP1\_UP\_EN}{1}{4}{{1}}%
\regfield{MCPWM\_OP0\_FORCE\_UP}{1}{3}{{0}}%
\regfield{MCPWM\_OP0\_UP\_EN}{1}{2}{{1}}%
\regfield{MCPWM\_GLOBAL\_FORCE\_UP}{1}{1}{{0}}%
\regfield{MCPWM\_GLOBAL\_UP\_EN}{1}{0}{{1}}%
\reglabel{Reset}\regnewline%
\begin{regdesc}\begin{reglist}
\label{fielddesc:MCPWMGLOBALUPEN}\item [MCPWM\_GLOBAL\_UP\_EN] 配置是否更新 MCPWM 模块的所有有效寄存器。 \\
0：无效\\
1：更新 MCPWM 模块的所有有效寄存器 \\
(R/W)
\label{fielddesc:MCPWMGLOBALFORCEUP}\item [MCPWM\_GLOBAL\_FORCE\_UP] 配置是否强制更新 MCPWM 模块的所有有效寄存器。\\
0：无效\\
1：强制更新 MCPWM 模块的所有有效寄存器 \\
(R/W)
\label{fielddesc:MCPWMOP0UPEN}\item [MCPWM\_OP0\_UP\_EN] 当 \hyperref[fielddesc:MCPWMGLOBALUPEN]{MCPWM\_GLOBAL\_UP\_EN} 置 1 时，配置是否更新 PWM 操作器 0 中的有效寄存器。 \\
0：无效\\
1：更新 PWM 操作器 0 中的有效寄存器 \\
(R/W)
\label{fielddesc:MCPWMOP0FORCEUP}\item [MCPWM\_OP0\_FORCE\_UP] 配置是否强制更新 PWM 操作器 0 中的有效寄存器。\\
0：无效\\
1：强制更新 \\
(R/W)
\label{fielddesc:MCPWMOP1UPEN}\item [MCPWM\_OP1\_UP\_EN] 当 \hyperref[fielddesc:MCPWMGLOBALUPEN]{MCPWM\_GLOBAL\_UP\_EN} 置 1 时，配置是否更新 PWM 操作器 1 中的有效寄存器。 \\
0：无效\\
1：更新 PWM 操作器 1 中的有效寄存器 \\
(R/W)
\label{fielddesc:MCPWMOP1FORCEUP}\item [MCPWM\_OP1\_FORCE\_UP] 配置是否强制更新 PWM 操作器 1 中的有效寄存器。\\
0：无效\\
1：强制更新 \\
(R/W)
\label{fielddesc:MCPWMOP2UPEN}\item [MCPWM\_OP2\_UP\_EN] 当 \hyperref[fielddesc:MCPWMGLOBALUPEN]{MCPWM\_GLOBAL\_UP\_EN} 置 1 时，配置是否更新 PWM 操作器 2 中的有效寄存器。 \\
0：无效\\
1：更新 PWM 操作器 2 中的有效寄存器 \\
(R/W)
\item [见下页...]
\end{reglist}\end{regdesc}
\end{register}
\addtocounter{Regfloat}{-1}
\begin{register}{H}{MCPWM\_UPDATE\_CFG\_REG}{0x010C}
\begin{regdesc}\begin{reglist}
\item [...接上页]
\label{fielddesc:MCPWMOP2FORCEUP}\item [MCPWM\_OP2\_FORCE\_UP] 配置是否强制更新 PWM 操作器 2 中的有效寄存器。\\
0：无效\\
1：强制更新 \\
(R/W)
\end{reglist}\end{regdesc}
\end{register}


\begin{register}{H}{MCPWM\_INT\_ENA\_REG}{0x0110}\label{regdesc:MCPWMINTENAREG}
\regfield{(reserved)}{2}{30}{00}%
\regfield{MCPWM\_CAP2\_INT\_ENA}{1}{29}{{0}}%
\regfield{MCPWM\_CAP1\_INT\_ENA}{1}{28}{{0}}%
\regfield{MCPWM\_CAP0\_INT\_ENA}{1}{27}{{0}}%
\regfield{MCPWM\_TZ2\_OST\_INT\_ENA}{1}{26}{{0}}%
\regfield{MCPWM\_TZ1\_OST\_INT\_ENA}{1}{25}{{0}}%
\regfield{MCPWM\_TZ0\_OST\_INT\_ENA}{1}{24}{{0}}%
\regfield{MCPWM\_TZ2\_CBC\_INT\_ENA}{1}{23}{{0}}%
\regfield{MCPWM\_TZ1\_CBC\_INT\_ENA}{1}{22}{{0}}%
\regfield{MCPWM\_TZ0\_CBC\_INT\_ENA}{1}{21}{{0}}%
\regfield{MCPWM\_CMPR2\_TEB\_INT\_ENA}{1}{20}{{0}}%
\regfield{MCPWM\_CMPR1\_TEB\_INT\_ENA}{1}{19}{{0}}%
\regfield{MCPWM\_CMPR0\_TEB\_INT\_ENA}{1}{18}{{0}}%
\regfield{MCPWM\_CMPR2\_TEA\_INT\_ENA}{1}{17}{{0}}%
\regfield{MCPWM\_CMPR1\_TEA\_INT\_ENA}{1}{16}{{0}}%
\regfield{MCPWM\_CMPR0\_TEA\_INT\_ENA}{1}{15}{{0}}%
\regfield{MCPWM\_FAULT2\_CLR\_INT\_ENA}{1}{14}{{0}}%
\regfield{MCPWM\_FAULT1\_CLR\_INT\_ENA}{1}{13}{{0}}%
\regfield{MCPWM\_FAULT0\_CLR\_INT\_ENA}{1}{12}{{0}}%
\regfield{MCPWM\_FAULT2\_INT\_ENA}{1}{11}{{0}}%
\regfield{MCPWM\_FAULT1\_INT\_ENA}{1}{10}{{0}}%
\regfield{MCPWM\_FAULT0\_INT\_ENA}{1}{9}{{0}}%
\regfield{MCPWM\_TIMER2\_TEP\_INT\_ENA}{1}{8}{{0}}%
\regfield{MCPWM\_TIMER1\_TEP\_INT\_ENA}{1}{7}{{0}}%
\regfield{MCPWM\_TIMER0\_TEP\_INT\_ENA}{1}{6}{{0}}%
\regfield{MCPWM\_TIMER2\_TEZ\_INT\_ENA}{1}{5}{{0}}%
\regfield{MCPWM\_TIMER1\_TEZ\_INT\_ENA}{1}{4}{{0}}%
\regfield{MCPWM\_TIMER0\_TEZ\_INT\_ENA}{1}{3}{{0}}%
\regfield{MCPWM\_TIMER2\_STOP\_INT\_ENA}{1}{2}{{0}}%
\regfield{MCPWM\_TIMER1\_STOP\_INT\_ENA}{1}{1}{{0}}%
\regfield{MCPWM\_TIMER0\_STOP\_INT\_ENA}{1}{0}{{0}}%
\reglabel{Reset}\regnewline%
\begin{regdesc}\begin{reglist}
\label{fielddesc:MCPWMTIMER0STOPINTENA}\item [MCPWM\_TIMER0\_STOP\_INT\_ENA] 写 1 使能 \hyperref[sec:mcpwm-interrupts]{MCPWM\_TIMER0\_STOP\_INT}。 (R/W)
\label{fielddesc:MCPWMTIMER1STOPINTENA}\item [MCPWM\_TIMER1\_STOP\_INT\_ENA] 写 1 使能 \hyperref[sec:mcpwm-interrupts]{MCPWM\_TIMER1\_STOP\_INT}。 (R/W)
\label{fielddesc:MCPWMTIMER2STOPINTENA}\item [MCPWM\_TIMER2\_STOP\_INT\_ENA] 写 1 使能 \hyperref[sec:mcpwm-interrupts]{MCPWM\_TIMER2\_STOP\_INT}。 (R/W)
\label{fielddesc:MCPWMTIMER0TEZINTENA}\item [MCPWM\_TIMER0\_TEZ\_INT\_ENA] 写 1 使能 \hyperref[sec:mcpwm-interrupts]{MCPWM\_TIMER0\_TEZ\_INT}。 (R/W)
\label{fielddesc:MCPWMTIMER1TEZINTENA}\item [MCPWM\_TIMER1\_TEZ\_INT\_ENA] 写 1 使能 \hyperref[sec:mcpwm-interrupts]{MCPWM\_TIMER1\_TEZ\_INT}。 (R/W)
\label{fielddesc:MCPWMTIMER2TEZINTENA}\item [MCPWM\_TIMER2\_TEZ\_INT\_ENA] 写 1 使能 \hyperref[sec:mcpwm-interrupts]{MCPWM\_TIMER2\_TEZ\_INT}。 (R/W)
\label{fielddesc:MCPWMTIMER0TEPINTENA}\item [MCPWM\_TIMER0\_TEP\_INT\_ENA] 写 1 使能 \hyperref[sec:mcpwm-interrupts]{MCPWM\_TIMER0\_TEP\_INT}。 (R/W)
\label{fielddesc:MCPWMTIMER1TEPINTENA}\item [MCPWM\_TIMER1\_TEP\_INT\_ENA] 写 1 使能 \hyperref[sec:mcpwm-interrupts]{MCPWM\_TIMER1\_TEP\_INT}。 (R/W)
\label{fielddesc:MCPWMTIMER2TEPINTENA}\item [MCPWM\_TIMER2\_TEP\_INT\_ENA] 写 1 使能 \hyperref[sec:mcpwm-interrupts]{MCPWM\_TIMER2\_TEP\_INT}。 (R/W)
\label{fielddesc:MCPWMFAULT0INTENA}\item [MCPWM\_FAULT0\_INT\_ENA] 写 1 使能 \hyperref[sec:mcpwm-interrupts]{MCPWM\_FAULT0\_INT}。 (R/W)
\label{fielddesc:MCPWMFAULT1INTENA}\item [MCPWM\_FAULT1\_INT\_ENA] 写 1 使能 \hyperref[sec:mcpwm-interrupts]{MCPWM\_FAULT1\_INT}。 (R/W)
\label{fielddesc:MCPWMFAULT2INTENA}\item [MCPWM\_FAULT2\_INT\_ENA] 写 1 使能 \hyperref[sec:mcpwm-interrupts]{MCPWM\_FAULT2\_INT}。 (R/W)
\label{fielddesc:MCPWMFAULT0CLRINTENA}\item [MCPWM\_FAULT0\_CLR\_INT\_ENA] 写 1 使能 \hyperref[sec:mcpwm-interrupts]{MCPWM\_FAULT0\_CLR\_INT}。 (R/W)
\label{fielddesc:MCPWMFAULT1CLRINTENA}\item [MCPWM\_FAULT1\_CLR\_INT\_ENA] 写 1 使能 \hyperref[sec:mcpwm-interrupts]{MCPWM\_FAULT1\_CLR\_INT}。 (R/W)
\label{fielddesc:MCPWMFAULT2CLRINTENA}\item [MCPWM\_FAULT2\_CLR\_INT\_ENA] 写 1 使能 \hyperref[sec:mcpwm-interrupts]{MCPWM\_FAULT2\_CLR\_INT}。(R/W)
\label{fielddesc:MCPWMCMPR0TEAINTENA}\item [MCPWM\_CMPR0\_TEA\_INT\_ENA] 写 1 使能 \hyperref[sec:mcpwm-interrupts]{MCPWM\_CMPR0\_TEA\_INT}。 (R/W)
\label{fielddesc:MCPWMCMPR1TEAINTENA}\item [MCPWM\_CMPR1\_TEA\_INT\_ENA] 写 1 使能 \hyperref[sec:mcpwm-interrupts]{MCPWM\_CMPR1\_TEA\_INT}。 (R/W)
\label{fielddesc:MCPWMCMPR2TEAINTENA}\item [MCPWM\_CMPR2\_TEA\_INT\_ENA] 写 1 使能 \hyperref[sec:mcpwm-interrupts]{MCPWM\_CMPR2\_TEA\_INT}。 (R/W)
\item [见下页...]
\end{reglist}\end{regdesc}
\end{register}
\addtocounter{Regfloat}{-1}
\begin{register}{H}{MCPWM\_INT\_ENA\_REG}{0x0110}
\begin{regdesc}\begin{reglist}
\item [...接上页]
\label{fielddesc:MCPWMCMPR0TEBINTENA}\item [MCPWM\_CMPR0\_TEB\_INT\_ENA] 写 1 使能 \hyperref[sec:mcpwm-interrupts]{MCPWM\_CMPR0\_TEB\_INT}。 (R/W)
\label{fielddesc:MCPWMCMPR1TEBINTENA}\item [MCPWM\_CMPR1\_TEB\_INT\_ENA] 写 1 使能 \hyperref[sec:mcpwm-interrupts]{MCPWM\_CMPR1\_TEB\_INT}。 (R/W)
\label{fielddesc:MCPWMCMPR2TEBINTENA}\item [MCPWM\_CMPR2\_TEB\_INT\_ENA] 写 1 使能 \hyperref[sec:mcpwm-interrupts]{MCPWM\_CMPR2\_TEB\_INT}。 (R/W)
\label{fielddesc:MCPWMTZ0CBCINTENA}\item [MCPWM\_TZ0\_CBC\_INT\_ENA] 写 1 使能 \hyperref[sec:mcpwm-interrupts]{MCPWM\_TZ0\_CBC\_INT}。 (R/W)
\label{fielddesc:MCPWMTZ1CBCINTENA}\item [MCPWM\_TZ1\_CBC\_INT\_ENA] 写 1 使能 \hyperref[sec:mcpwm-interrupts]{MCPWM\_TZ1\_CBC\_INT}。 (R/W)
\label{fielddesc:MCPWMTZ2CBCINTENA}\item [MCPWM\_TZ2\_CBC\_INT\_ENA] 写 1 使能 \hyperref[sec:mcpwm-interrupts]{MCPWM\_TZ2\_CBC\_INT}。 (R/W)
\label{fielddesc:MCPWMTZ0OSTINTENA}\item [MCPWM\_TZ0\_OST\_INT\_ENA] 写 1 使能 \hyperref[sec:mcpwm-interrupts]{MCPWM\_TZ0\_OST\_INT}。 (R/W)
\label{fielddesc:MCPWMTZ1OSTINTENA}\item [MCPWM\_TZ1\_OST\_INT\_ENA] 写 1 使能 \hyperref[sec:mcpwm-interrupts]{MCPWM\_TZ1\_OST\_INT}。 (R/W)
\label{fielddesc:MCPWMTZ2OSTINTENA}\item [MCPWM\_TZ2\_OST\_INT\_ENA] 写 1 使能 \hyperref[sec:mcpwm-interrupts]{MCPWM\_TZ2\_OST\_INT}。 (R/W)
\label{fielddesc:MCPWMCAP0INTENA}\item [MCPWM\_CAP0\_INT\_ENA] 写 1 使能 \hyperref[sec:mcpwm-interrupts]{MCPWM\_CAP0\_INT}。 (R/W)
\label{fielddesc:MCPWMCAP1INTENA}\item [MCPWM\_CAP1\_INT\_ENA] 写 1 使能 \hyperref[sec:mcpwm-interrupts]{MCPWM\_CAP1\_INT}。 (R/W)
\label{fielddesc:MCPWMCAP2INTENA}\item [MCPWM\_CAP2\_INT\_ENA] 写 1 使能 \hyperref[sec:mcpwm-interrupts]{MCPWM\_CAP2\_INT}。 (R/W)
\end{reglist}\end{regdesc}
\end{register}


\begin{register}{H}{MCPWM\_INT\_RAW\_REG}{0x0114}\label{regdesc:MCPWMINTRAWREG}
\regfield{(reserved)}{2}{30}{00}%
\regfield{MCPWM\_CAP2\_INT\_RAW}{1}{29}{{0}}%
\regfield{MCPWM\_CAP1\_INT\_RAW}{1}{28}{{0}}%
\regfield{MCPWM\_CAP0\_INT\_RAW}{1}{27}{{0}}%
\regfield{MCPWM\_TZ2\_OST\_INT\_RAW}{1}{26}{{0}}%
\regfield{MCPWM\_TZ1\_OST\_INT\_RAW}{1}{25}{{0}}%
\regfield{MCPWM\_TZ0\_OST\_INT\_RAW}{1}{24}{{0}}%
\regfield{MCPWM\_TZ2\_CBC\_INT\_RAW}{1}{23}{{0}}%
\regfield{MCPWM\_TZ1\_CBC\_INT\_RAW}{1}{22}{{0}}%
\regfield{MCPWM\_TZ0\_CBC\_INT\_RAW}{1}{21}{{0}}%
\regfield{MCPWM\_CMPR2\_TEB\_INT\_RAW}{1}{20}{{0}}%
\regfield{MCPWM\_CMPR1\_TEB\_INT\_RAW}{1}{19}{{0}}%
\regfield{MCPWM\_CMPR0\_TEB\_INT\_RAW}{1}{18}{{0}}%
\regfield{MCPWM\_CMPR2\_TEA\_INT\_RAW}{1}{17}{{0}}%
\regfield{MCPWM\_CMPR1\_TEA\_INT\_RAW}{1}{16}{{0}}%
\regfield{MCPWM\_CMPR0\_TEA\_INT\_RAW}{1}{15}{{0}}%
\regfield{MCPWM\_FAULT2\_CLR\_INT\_RAW}{1}{14}{{0}}%
\regfield{MCPWM\_FAULT1\_CLR\_INT\_RAW}{1}{13}{{0}}%
\regfield{MCPWM\_FAULT0\_CLR\_INT\_RAW}{1}{12}{{0}}%
\regfield{MCPWM\_FAULT2\_INT\_RAW}{1}{11}{{0}}%
\regfield{MCPWM\_FAULT1\_INT\_RAW}{1}{10}{{0}}%
\regfield{MCPWM\_FAULT0\_INT\_RAW}{1}{9}{{0}}%
\regfield{MCPWM\_TIMER2\_TEP\_INT\_RAW}{1}{8}{{0}}%
\regfield{MCPWM\_TIMER1\_TEP\_INT\_RAW}{1}{7}{{0}}%
\regfield{MCPWM\_TIMER0\_TEP\_INT\_RAW}{1}{6}{{0}}%
\regfield{MCPWM\_TIMER2\_TEZ\_INT\_RAW}{1}{5}{{0}}%
\regfield{MCPWM\_TIMER1\_TEZ\_INT\_RAW}{1}{4}{{0}}%
\regfield{MCPWM\_TIMER0\_TEZ\_INT\_RAW}{1}{3}{{0}}%
\regfield{MCPWM\_TIMER2\_STOP\_INT\_RAW}{1}{2}{{0}}%
\regfield{MCPWM\_TIMER1\_STOP\_INT\_RAW}{1}{1}{{0}}%
\regfield{MCPWM\_TIMER0\_STOP\_INT\_RAW}{1}{0}{{0}}%
\reglabel{Reset}\regnewline%
\begin{regdesc}\begin{reglist}
\label{fielddesc:MCPWMTIMER0STOPINTRAW}\item [MCPWM\_TIMER0\_STOP\_INT\_RAW] \hyperref[sec:mcpwm-interrupts]{MCPWM\_TIMER0\_STOP\_INT} 的原始状态位。 (R/WTC/SS)
\label{fielddesc:MCPWMTIMER1STOPINTRAW}\item [MCPWM\_TIMER1\_STOP\_INT\_RAW] \hyperref[sec:mcpwm-interrupts]{MCPWM\_TIMER1\_STOP\_INT} 的原始状态位。 (R/WTC/SS)
\label{fielddesc:MCPWMTIMER2STOPINTRAW}\item [MCPWM\_TIMER2\_STOP\_INT\_RAW] \hyperref[sec:mcpwm-interrupts]{MCPWM\_TIMER2\_STOP\_INT} 的原始状态位。 (R/WTC/SS)
\label{fielddesc:MCPWMTIMER0TEZINTRAW}\item [MCPWM\_TIMER0\_TEZ\_INT\_RAW] \hyperref[sec:mcpwm-interrupts]{MCPWM\_TIMER0\_TEZ\_INT} 的原始状态位。 (R/WTC/SS)
\label{fielddesc:MCPWMTIMER1TEZINTRAW}\item [MCPWM\_TIMER1\_TEZ\_INT\_RAW] \hyperref[sec:mcpwm-interrupts]{MCPWM\_TIMER1\_TEZ\_INT} 的原始状态位。 (R/WTC/SS)
\label{fielddesc:MCPWMTIMER2TEZINTRAW}\item [MCPWM\_TIMER2\_TEZ\_INT\_RAW] \hyperref[sec:mcpwm-interrupts]{MCPWM\_TIMER2\_TEZ\_INT} 的原始状态位。 (R/WTC/SS)
\label{fielddesc:MCPWMTIMER0TEPINTRAW}\item [MCPWM\_TIMER0\_TEP\_INT\_RAW] \hyperref[sec:mcpwm-interrupts]{MCPWM\_TIMER0\_TEP\_INT} 的原始状态位。 (R/WTC/SS)
\label{fielddesc:MCPWMTIMER1TEPINTRAW}\item [MCPWM\_TIMER1\_TEP\_INT\_RAW] \hyperref[sec:mcpwm-interrupts]{MCPWM\_TIMER1\_TEP\_INT} 的原始状态位。 (R/WTC/SS)
\label{fielddesc:MCPWMTIMER2TEPINTRAW}\item [MCPWM\_TIMER2\_TEP\_INT\_RAW] \hyperref[sec:mcpwm-interrupts]{MCPWM\_TIMER2\_TEP\_INT} 的原始状态位。 (R/WTC/SS)
\label{fielddesc:MCPWMFAULT0INTRAW}\item [MCPWM\_FAULT0\_INT\_RAW] \hyperref[sec:mcpwm-interrupts]{MCPWM\_FAULT0\_INT} 的原始状态位。 (R/WTC/SS)
\label{fielddesc:MCPWMFAULT1INTRAW}\item [MCPWM\_FAULT1\_INT\_RAW] \hyperref[sec:mcpwm-interrupts]{MCPWM\_FAULT1\_INT} 的原始状态位。 (R/WTC/SS)
\label{fielddesc:MCPWMFAULT2INTRAW}\item [MCPWM\_FAULT2\_INT\_RAW] \hyperref[sec:mcpwm-interrupts]{MCPWM\_FAULT2\_INT} 的原始状态位。 (R/WTC/SS)
\label{fielddesc:MCPWMFAULT0CLRINTRAW}\item [MCPWM\_FAULT0\_CLR\_INT\_RAW] \hyperref[sec:mcpwm-interrupts]{MCPWM\_FAULT0\_CLR\_INT} 的原始状态位。 (R/WTC/SS)
\label{fielddesc:MCPWMFAULT1CLRINTRAW}\item [MCPWM\_FAULT1\_CLR\_INT\_RAW] \hyperref[sec:mcpwm-interrupts]{MCPWM\_FAULT1\_CLR\_INT} 的原始状态位。 (R/WTC/SS)
\label{fielddesc:MCPWMFAULT2CLRINTRAW}\item [MCPWM\_FAULT2\_CLR\_INT\_RAW] \hyperref[sec:mcpwm-interrupts]{MCPWM\_FAULT2\_CLR\_INT} 的原始状态位。 (R/WTC/SS)
\item [见下页...]
\end{reglist}\end{regdesc}
\end{register}
\addtocounter{Regfloat}{-1}
\begin{register}{H}{MCPWM\_INT\_RAW\_REG}{0x0114}
\begin{regdesc}\begin{reglist}
\item [...接上页]
\label{fielddesc:MCPWMCMPR0TEAINTRAW}\item [MCPWM\_CMPR0\_TEA\_INT\_RAW] \hyperref[sec:mcpwm-interrupts]{MCPWM\_CMPR0\_TEA\_INT} 的原始状态位。 (R/WTC/SS)
\label{fielddesc:MCPWMCMPR1TEAINTRAW}\item [MCPWM\_CMPR1\_TEA\_INT\_RAW] \hyperref[sec:mcpwm-interrupts]{MCPWM\_CMPR1\_TEA\_INT} 的原始状态位。 (R/WTC/SS)
\label{fielddesc:MCPWMCMPR2TEAINTRAW}\item [MCPWM\_CMPR2\_TEA\_INT\_RAW] \hyperref[sec:mcpwm-interrupts]{MCPWM\_CMPR2\_TEA\_INT} 的原始状态位。 (R/WTC/SS)
\label{fielddesc:MCPWMCMPR0TEBINTRAW}\item [MCPWM\_CMPR0\_TEB\_INT\_RAW] \hyperref[sec:mcpwm-interrupts]{MCPWM\_CMPR0\_TEB\_INT} 的原始状态位。 (R/WTC/SS)
\label{fielddesc:MCPWMCMPR1TEBINTRAW}\item [MCPWM\_CMPR1\_TEB\_INT\_RAW] \hyperref[sec:mcpwm-interrupts]{MCPWM\_CMPR1\_TEB\_INT} 的原始状态位。 (R/WTC/SS)
\label{fielddesc:MCPWMCMPR2TEBINTRAW}\item [MCPWM\_CMPR2\_TEB\_INT\_RAW] \hyperref[sec:mcpwm-interrupts]{MCPWM\_CMPR2\_TEB\_INT} 的原始状态位。 (R/WTC/SS)
\label{fielddesc:MCPWMTZ0CBCINTRAW}\item [MCPWM\_TZ0\_CBC\_INT\_RAW] \hyperref[sec:mcpwm-interrupts]{MCPWM\_TZ0\_CBC\_INT} 的原始状态位。 (R/WTC/SS)
\label{fielddesc:MCPWMTZ1CBCINTRAW}\item [MCPWM\_TZ1\_CBC\_INT\_RAW] \hyperref[sec:mcpwm-interrupts]{MCPWM\_TZ1\_CBC\_INT} 的原始状态位。 (R/WTC/SS)
\label{fielddesc:MCPWMTZ2CBCINTRAW}\item [MCPWM\_TZ2\_CBC\_INT\_RAW] \hyperref[sec:mcpwm-interrupts]{MCPWM\_TZ2\_CBC\_INT} 的原始状态位。 (R/WTC/SS)
\label{fielddesc:MCPWMTZ0OSTINTRAW}\item [MCPWM\_TZ0\_OST\_INT\_RAW] \hyperref[sec:mcpwm-interrupts]{MCPWM\_TZ0\_OST\_INT} 的原始状态位。 (R/WTC/SS)
\label{fielddesc:MCPWMTZ1OSTINTRAW}\item [MCPWM\_TZ1\_OST\_INT\_RAW] \hyperref[sec:mcpwm-interrupts]{MCPWM\_TZ1\_OST\_INT} 的原始状态位。 (R/WTC/SS)
\label{fielddesc:MCPWMTZ2OSTINTRAW}\item [MCPWM\_TZ2\_OST\_INT\_RAW] \hyperref[sec:mcpwm-interrupts]{MCPWM\_TZ2\_OST\_INT} 的原始状态位。 (R/WTC/SS)
\label{fielddesc:MCPWMCAP0INTRAW}\item [MCPWM\_CAP0\_INT\_RAW] \hyperref[sec:mcpwm-interrupts]{MCPWM\_CAP0\_INT} 的原始状态位。 (R/WTC/SS)
\label{fielddesc:MCPWMCAP1INTRAW}\item [MCPWM\_CAP1\_INT\_RAW] \hyperref[sec:mcpwm-interrupts]{MCPWM\_CAP1\_INT} 的原始状态位。 (R/WTC/SS)
\label{fielddesc:MCPWMCAP2INTRAW}\item [MCPWM\_CAP2\_INT\_RAW] \hyperref[sec:mcpwm-interrupts]{MCPWM\_CAP2\_INT} 的原始状态位。 (R/WTC/SS)
\end{reglist}\end{regdesc}
\end{register}


\begin{register}{H}{MCPWM\_INT\_ST\_REG}{0x0118}\label{regdesc:MCPWMINTSTREG}
\regfield{(reserved)}{2}{30}{00}%
\regfield{MCPWM\_CAP2\_INT\_ST}{1}{29}{{0}}%
\regfield{MCPWM\_CAP1\_INT\_ST}{1}{28}{{0}}%
\regfield{MCPWM\_CAP0\_INT\_ST}{1}{27}{{0}}%
\regfield{MCPWM\_TZ2\_OST\_INT\_ST}{1}{26}{{0}}%
\regfield{MCPWM\_TZ1\_OST\_INT\_ST}{1}{25}{{0}}%
\regfield{MCPWM\_TZ0\_OST\_INT\_ST}{1}{24}{{0}}%
\regfield{MCPWM\_TZ2\_CBC\_INT\_ST}{1}{23}{{0}}%
\regfield{MCPWM\_TZ1\_CBC\_INT\_ST}{1}{22}{{0}}%
\regfield{MCPWM\_TZ0\_CBC\_INT\_ST}{1}{21}{{0}}%
\regfield{MCPWM\_CMPR2\_TEB\_INT\_ST}{1}{20}{{0}}%
\regfield{MCPWM\_CMPR1\_TEB\_INT\_ST}{1}{19}{{0}}%
\regfield{MCPWM\_CMPR0\_TEB\_INT\_ST}{1}{18}{{0}}%
\regfield{MCPWM\_CMPR2\_TEA\_INT\_ST}{1}{17}{{0}}%
\regfield{MCPWM\_CMPR1\_TEA\_INT\_ST}{1}{16}{{0}}%
\regfield{MCPWM\_CMPR0\_TEA\_INT\_ST}{1}{15}{{0}}%
\regfield{MCPWM\_FAULT2\_CLR\_INT\_ST}{1}{14}{{0}}%
\regfield{MCPWM\_FAULT1\_CLR\_INT\_ST}{1}{13}{{0}}%
\regfield{MCPWM\_FAULT0\_CLR\_INT\_ST}{1}{12}{{0}}%
\regfield{MCPWM\_FAULT2\_INT\_ST}{1}{11}{{0}}%
\regfield{MCPWM\_FAULT1\_INT\_ST}{1}{10}{{0}}%
\regfield{MCPWM\_FAULT0\_INT\_ST}{1}{9}{{0}}%
\regfield{MCPWM\_TIMER2\_TEP\_INT\_ST}{1}{8}{{0}}%
\regfield{MCPWM\_TIMER1\_TEP\_INT\_ST}{1}{7}{{0}}%
\regfield{MCPWM\_TIMER0\_TEP\_INT\_ST}{1}{6}{{0}}%
\regfield{MCPWM\_TIMER2\_TEZ\_INT\_ST}{1}{5}{{0}}%
\regfield{MCPWM\_TIMER1\_TEZ\_INT\_ST}{1}{4}{{0}}%
\regfield{MCPWM\_TIMER0\_TEZ\_INT\_ST}{1}{3}{{0}}%
\regfield{MCPWM\_TIMER2\_STOP\_INT\_ST}{1}{2}{{0}}%
\regfield{MCPWM\_TIMER1\_STOP\_INT\_ST}{1}{1}{{0}}%
\regfield{MCPWM\_TIMER0\_STOP\_INT\_ST}{1}{0}{{0}}%
\reglabel{Reset}\regnewline%
\begin{regdesc}\begin{reglist}
\label{fielddesc:MCPWMTIMER0STOPINTST}\item [MCPWM\_TIMER0\_STOP\_INT\_ST] \hyperref[sec:mcpwm-interrupts]{MCPWM\_TIMER0\_STOP\_INT} 的屏蔽状态位。 (RO)
\label{fielddesc:MCPWMTIMER1STOPINTST}\item [MCPWM\_TIMER1\_STOP\_INT\_ST] \hyperref[sec:mcpwm-interrupts]{MCPWM\_TIMER1\_STOP\_INT} 的屏蔽状态位。 (RO)
\label{fielddesc:MCPWMTIMER2STOPINTST}\item [MCPWM\_TIMER2\_STOP\_INT\_ST] \hyperref[sec:mcpwm-interrupts]{MCPWM\_TIMER2\_STOP\_INT} 的屏蔽状态位。 (RO)
\label{fielddesc:MCPWMTIMER0TEZINTST}\item [MCPWM\_TIMER0\_TEZ\_INT\_ST] \hyperref[sec:mcpwm-interrupts]{MCPWM\_TIMER0\_TEZ\_INT} 的屏蔽状态位。 (RO)
\label{fielddesc:MCPWMTIMER1TEZINTST}\item [MCPWM\_TIMER1\_TEZ\_INT\_ST] \hyperref[sec:mcpwm-interrupts]{MCPWM\_TIMER1\_TEZ\_INT} 的屏蔽状态位。 (RO)
\label{fielddesc:MCPWMTIMER2TEZINTST}\item [MCPWM\_TIMER2\_TEZ\_INT\_ST] \hyperref[sec:mcpwm-interrupts]{MCPWM\_TIMER2\_TEZ\_INT} 的屏蔽状态位。 (RO)
\label{fielddesc:MCPWMTIMER0TEPINTST}\item [MCPWM\_TIMER0\_TEP\_INT\_ST] \hyperref[sec:mcpwm-interrupts]{MCPWM\_TIMER0\_TEP\_INT} 的屏蔽状态位。 (RO)
\label{fielddesc:MCPWMTIMER1TEPINTST}\item [MCPWM\_TIMER1\_TEP\_INT\_ST] \hyperref[sec:mcpwm-interrupts]{MCPWM\_TIMER1\_TEP\_INT} 的屏蔽状态位。 (RO)
\label{fielddesc:MCPWMTIMER2TEPINTST}\item [MCPWM\_TIMER2\_TEP\_INT\_ST] \hyperref[sec:mcpwm-interrupts]{MCPWM\_TIMER2\_TEP\_INT} 的屏蔽状态位。 (RO)
\label{fielddesc:MCPWMFAULT0INTST}\item [MCPWM\_FAULT0\_INT\_ST] \hyperref[sec:mcpwm-interrupts]{MCPWM\_FAULT0\_INT\_INT} 的屏蔽状态位。 (RO)
\label{fielddesc:MCPWMFAULT1INTST}\item [MCPWM\_FAULT1\_INT\_ST] \hyperref[sec:mcpwm-interrupts]{MCPWM\_FAULT1\_INT\_INT} 的屏蔽状态位。 (RO)
\label{fielddesc:MCPWMFAULT2INTST}\item [MCPWM\_FAULT2\_INT\_ST] \hyperref[sec:mcpwm-interrupts]{MCPWM\_FAULT2\_INT\_INT} 的屏蔽状态位。 (RO)
\label{fielddesc:MCPWMFAULT0CLRINTST}\item [MCPWM\_FAULT0\_CLR\_INT\_ST] \hyperref[sec:mcpwm-interrupts]{MCPWM\_FAULT0\_CLR\_INT} 的屏蔽状态位。 (RO)
\label{fielddesc:MCPWMFAULT1CLRINTST}\item [MCPWM\_FAULT1\_CLR\_INT\_ST] \hyperref[sec:mcpwm-interrupts]{MCPWM\_FAULT1\_CLR\_INT} 的屏蔽状态位。 (RO)
\label{fielddesc:MCPWMFAULT2CLRINTST}\item [MCPWM\_FAULT2\_CLR\_INT\_ST] \hyperref[sec:mcpwm-interrupts]{MCPWM\_FAULT2\_CLR\_INT} 的屏蔽状态位。 (RO)
\label{fielddesc:MCPWMCMPR0TEAINTST}\item [MCPWM\_CMPR0\_TEA\_INT\_ST] \hyperref[sec:mcpwm-interrupts]{MCPWM\_CMPR0\_TEA\_INT} 的屏蔽状态位。 (RO)
\label{fielddesc:MCPWMCMPR1TEAINTST}\item [MCPWM\_CMPR1\_TEA\_INT\_ST] \hyperref[sec:mcpwm-interrupts]{MCPWM\_CMPR1\_TEA\_INT} 的屏蔽状态位。 (RO)
\label{fielddesc:MCPWMCMPR2TEAINTST}\item [MCPWM\_CMPR2\_TEA\_INT\_ST] \hyperref[sec:mcpwm-interrupts]{MCPWM\_CMPR2\_TEA\_INT} 的屏蔽状态位。 (RO)
\item [见下页...]
\end{reglist}\end{regdesc}
\end{register}
\addtocounter{Regfloat}{-1}
\begin{register}{H}{MCPWM\_INT\_ST\_REG}{0x0118}
\begin{regdesc}\begin{reglist}
\item [...接上页]
\label{fielddesc:MCPWMCMPR0TEBINTST}\item [MCPWM\_CMPR0\_TEB\_INT\_ST] \hyperref[sec:mcpwm-interrupts]{MCPWM\_CMPR0\_TEB\_INT} 的屏蔽状态位。 (RO)
\label{fielddesc:MCPWMCMPR1TEBINTST}\item [MCPWM\_CMPR1\_TEB\_INT\_ST] \hyperref[sec:mcpwm-interrupts]{MCPWM\_CMPR1\_TEB\_INT} 的屏蔽状态位。 (RO)
\label{fielddesc:MCPWMCMPR2TEBINTST}\item [MCPWM\_CMPR2\_TEB\_INT\_ST] \hyperref[sec:mcpwm-interrupts]{MCPWM\_CMPR2\_TEB\_INT} 的屏蔽状态位。 (RO)
\label{fielddesc:MCPWMTZ0CBCINTST}\item [MCPWM\_TZ0\_CBC\_INT\_ST] \hyperref[sec:mcpwm-interrupts]{MCPWM\_TZ0\_CBC\_INT} 的屏蔽状态位。 (RO)
\label{fielddesc:MCPWMTZ1CBCINTST}\item [MCPWM\_TZ1\_CBC\_INT\_ST] \hyperref[sec:mcpwm-interrupts]{MCPWM\_TZ1\_CBC\_INT} 的屏蔽状态位。 (RO)
\label{fielddesc:MCPWMTZ2CBCINTST}\item [MCPWM\_TZ2\_CBC\_INT\_ST] \hyperref[sec:mcpwm-interrupts]{MCPWM\_TZ2\_CBC\_INT} 的屏蔽状态位。 (RO)
\label{fielddesc:MCPWMTZ0OSTINTST}\item [MCPWM\_TZ0\_OST\_INT\_ST] \hyperref[sec:mcpwm-interrupts]{MCPWM\_TZ0\_OST\_INT} 的屏蔽状态位。 (RO)
\label{fielddesc:MCPWMTZ1OSTINTST}\item [MCPWM\_TZ1\_OST\_INT\_ST] \hyperref[sec:mcpwm-interrupts]{MCPWM\_TZ1\_OST\_INT} 的屏蔽状态位。 (RO)
\label{fielddesc:MCPWMTZ2OSTINTST}\item [MCPWM\_TZ2\_OST\_INT\_ST] \hyperref[sec:mcpwm-interrupts]{MCPWM\_TZ2\_OST\_INT} 的屏蔽状态位。 (RO)
\label{fielddesc:MCPWMCAP0INTST}\item [MCPWM\_CAP0\_INT\_ST] \hyperref[sec:mcpwm-interrupts]{MCPWM\_CAP0\_INT} 的屏蔽状态位。 (RO)
\label{fielddesc:MCPWMCAP1INTST}\item [MCPWM\_CAP1\_INT\_ST] \hyperref[sec:mcpwm-interrupts]{MCPWM\_CAP1\_INT} 的屏蔽状态位。 (RO)
\label{fielddesc:MCPWMCAP2INTST}\item [MCPWM\_CAP2\_INT\_ST] \hyperref[sec:mcpwm-interrupts]{MCPWM\_CAP2\_INT} 的屏蔽状态位。 (RO)
\end{reglist}\end{regdesc}
\end{register}


\begin{register}{H}{MCPWM\_INT\_CLR\_REG}{0x011C}\label{regdesc:MCPWMINTCLRREG}
\regfield{(reserved)}{2}{30}{00}%
\regfield{MCPWM\_CAP2\_INT\_CLR}{1}{29}{{0}}%
\regfield{MCPWM\_CAP1\_INT\_CLR}{1}{28}{{0}}%
\regfield{MCPWM\_CAP0\_INT\_CLR}{1}{27}{{0}}%
\regfield{MCPWM\_TZ2\_OST\_INT\_CLR}{1}{26}{{0}}%
\regfield{MCPWM\_TZ1\_OST\_INT\_CLR}{1}{25}{{0}}%
\regfield{MCPWM\_TZ0\_OST\_INT\_CLR}{1}{24}{{0}}%
\regfield{MCPWM\_TZ2\_CBC\_INT\_CLR}{1}{23}{{0}}%
\regfield{MCPWM\_TZ1\_CBC\_INT\_CLR}{1}{22}{{0}}%
\regfield{MCPWM\_TZ0\_CBC\_INT\_CLR}{1}{21}{{0}}%
\regfield{MCPWM\_CMPR2\_TEB\_INT\_CLR}{1}{20}{{0}}%
\regfield{MCPWM\_CMPR1\_TEB\_INT\_CLR}{1}{19}{{0}}%
\regfield{MCPWM\_CMPR0\_TEB\_INT\_CLR}{1}{18}{{0}}%
\regfield{MCPWM\_CMPR2\_TEA\_INT\_CLR}{1}{17}{{0}}%
\regfield{MCPWM\_CMPR1\_TEA\_INT\_CLR}{1}{16}{{0}}%
\regfield{MCPWM\_CMPR0\_TEA\_INT\_CLR}{1}{15}{{0}}%
\regfield{MCPWM\_FAULT2\_CLR\_INT\_CLR}{1}{14}{{0}}%
\regfield{MCPWM\_FAULT1\_CLR\_INT\_CLR}{1}{13}{{0}}%
\regfield{MCPWM\_FAULT0\_CLR\_INT\_CLR}{1}{12}{{0}}%
\regfield{MCPWM\_FAULT2\_INT\_CLR}{1}{11}{{0}}%
\regfield{MCPWM\_FAULT1\_INT\_CLR}{1}{10}{{0}}%
\regfield{MCPWM\_FAULT0\_INT\_CLR}{1}{9}{{0}}%
\regfield{MCPWM\_TIMER2\_TEP\_INT\_CLR}{1}{8}{{0}}%
\regfield{MCPWM\_TIMER1\_TEP\_INT\_CLR}{1}{7}{{0}}%
\regfield{MCPWM\_TIMER0\_TEP\_INT\_CLR}{1}{6}{{0}}%
\regfield{MCPWM\_TIMER2\_TEZ\_INT\_CLR}{1}{5}{{0}}%
\regfield{MCPWM\_TIMER1\_TEZ\_INT\_CLR}{1}{4}{{0}}%
\regfield{MCPWM\_TIMER0\_TEZ\_INT\_CLR}{1}{3}{{0}}%
\regfield{MCPWM\_TIMER2\_STOP\_INT\_CLR}{1}{2}{{0}}%
\regfield{MCPWM\_TIMER1\_STOP\_INT\_CLR}{1}{1}{{0}}%
\regfield{MCPWM\_TIMER0\_STOP\_INT\_CLR}{1}{0}{{0}}%
\reglabel{Reset}\regnewline%
\begin{regdesc}\begin{reglist}
\label{fielddesc:MCPWMTIMER0STOPINTCLR}\item [MCPWM\_TIMER0\_STOP\_INT\_CLR] 写 1 清除 \hyperref[sec:mcpwm-interrupts]{MCPWM\_TIMER0\_STOP\_INT}。 (WT)
\label{fielddesc:MCPWMTIMER1STOPINTCLR}\item [MCPWM\_TIMER1\_STOP\_INT\_CLR] 写 1 清除 \hyperref[sec:mcpwm-interrupts]{MCPWM\_TIMER1\_STOP\_INT}。 (WT)
\label{fielddesc:MCPWMTIMER2STOPINTCLR}\item [MCPWM\_TIMER2\_STOP\_INT\_CLR] 写 1 清除 \hyperref[sec:mcpwm-interrupts]{MCPWM\_TIMER2\_STOP\_INT}。 (WT)
\label{fielddesc:MCPWMTIMER0TEZINTCLR}\item [MCPWM\_TIMER0\_TEZ\_INT\_CLR] 写 1 清除 \hyperref[sec:mcpwm-interrupts]{MCPWM\_TIMER0\_TEZ\_INT}。 (WT)
\label{fielddesc:MCPWMTIMER1TEZINTCLR}\item [MCPWM\_TIMER1\_TEZ\_INT\_CLR] 写 1 清除 \hyperref[sec:mcpwm-interrupts]{MCPWM\_TIMER1\_TEZ\_INT}。 (WT)
\label{fielddesc:MCPWMTIMER2TEZINTCLR}\item [MCPWM\_TIMER2\_TEZ\_INT\_CLR] 写 1 清除 \hyperref[sec:mcpwm-interrupts]{MCPWM\_TIMER2\_TEZ\_INT}。 (WT)
\label{fielddesc:MCPWMTIMER0TEPINTCLR}\item [MCPWM\_TIMER0\_TEP\_INT\_CLR] 写 1 清除 \hyperref[sec:mcpwm-interrupts]{MCPWM\_TIMER0\_TEP\_INT}。 (WT)
\label{fielddesc:MCPWMTIMER1TEPINTCLR}\item [MCPWM\_TIMER1\_TEP\_INT\_CLR] 写 1 清除 \hyperref[sec:mcpwm-interrupts]{MCPWM\_TIMER1\_TEP\_INT}。 (WT)
\label{fielddesc:MCPWMTIMER2TEPINTCLR}\item [MCPWM\_TIMER2\_TEP\_INT\_CLR] 写 1 清除 \hyperref[sec:mcpwm-interrupts]{MCPWM\_TIMER2\_TEP\_INT}。 (WT)
\label{fielddesc:MCPWMFAULT0INTCLR}\item [MCPWM\_FAULT0\_INT\_CLR] 写 1 清除 \hyperref[sec:mcpwm-interrupts]{MCPWM\_FAULT0\_INT}。 (WT)
\label{fielddesc:MCPWMFAULT1INTCLR}\item [MCPWM\_FAULT1\_INT\_CLR] 写 1 清除 \hyperref[sec:mcpwm-interrupts]{MCPWM\_FAULT1\_INT}。 (WT)
\label{fielddesc:MCPWMFAULT2INTCLR}\item [MCPWM\_FAULT2\_INT\_CLR] 写 1 清除 \hyperref[sec:mcpwm-interrupts]{MCPWM\_FAULT2\_INT}。 (WT)
\label{fielddesc:MCPWMFAULT0CLRINTCLR}\item [MCPWM\_FAULT0\_CLR\_INT\_CLR] 写 1 清除 \hyperref[sec:mcpwm-interrupts]{MCPWM\_FAULT0\_CLR\_INT}。 (WT)
\label{fielddesc:MCPWMFAULT1CLRINTCLR}\item [MCPWM\_FAULT1\_CLR\_INT\_CLR] 写 1 清除 \hyperref[sec:mcpwm-interrupts]{MCPWM\_FAULT1\_CLR\_INT}。 (WT)
\label{fielddesc:MCPWMFAULT2CLRINTCLR}\item [MCPWM\_FAULT2\_CLR\_INT\_CLR] 写 1 清除 \hyperref[sec:mcpwm-interrupts]{MCPWM\_FAULT2\_CLR\_INT}。 (WT)
\item [见下页...]
\end{reglist}\end{regdesc}
\end{register}
\addtocounter{Regfloat}{-1}
\begin{register}{H}{MCPWM\_INT\_CLR\_REG}{0x011C}
\begin{regdesc}\begin{reglist}
\item [...接上页]
\label{fielddesc:MCPWMCMPR0TEAINTCLR}\item [MCPWM\_CMPR0\_TEA\_INT\_CLR] 写 1 清除 \hyperref[sec:mcpwm-interrupts]{MCPWM\_CMPR0\_TEA\_INT}。 (WT)
\label{fielddesc:MCPWMCMPR1TEAINTCLR}\item [MCPWM\_CMPR1\_TEA\_INT\_CLR] 写 1 清除 \hyperref[sec:mcpwm-interrupts]{MCPWM\_CMPR1\_TEA\_INT}。 (WT)
\label{fielddesc:MCPWMCMPR2TEAINTCLR}\item [MCPWM\_CMPR2\_TEA\_INT\_CLR] 写 1 清除 \hyperref[sec:mcpwm-interrupts]{MCPWM\_CMPR2\_TEA\_INT}。 (WT)
\label{fielddesc:MCPWMCMPR0TEBINTCLR}\item [MCPWM\_CMPR0\_TEB\_INT\_CLR] 写 1 清除 \hyperref[sec:mcpwm-interrupts]{MCPWM\_CMPR0\_TEB\_INT}。 (WT)
\label{fielddesc:MCPWMCMPR1TEBINTCLR}\item [MCPWM\_CMPR1\_TEB\_INT\_CLR] 写 1 清除 \hyperref[sec:mcpwm-interrupts]{MCPWM\_CMPR1\_TEB\_INT}。 (WT)
\label{fielddesc:MCPWMCMPR2TEBINTCLR}\item [MCPWM\_CMPR2\_TEB\_INT\_CLR] 写 1 清除 \hyperref[sec:mcpwm-interrupts]{MCPWM\_CMPR2\_TEB\_INT}。 (WT)
\label{fielddesc:MCPWMTZ0CBCINTCLR}\item [MCPWM\_TZ0\_CBC\_INT\_CLR] 写 1 清除 \hyperref[sec:mcpwm-interrupts]{MCPWM\_TZ0\_CBC\_INT}。 (WT)
\label{fielddesc:MCPWMTZ1CBCINTCLR}\item [MCPWM\_TZ1\_CBC\_INT\_CLR] 写 1 清除 \hyperref[sec:mcpwm-interrupts]{MCPWM\_TZ1\_CBC\_INT}。 (WT)
\label{fielddesc:MCPWMTZ2CBCINTCLR}\item [MCPWM\_TZ2\_CBC\_INT\_CLR] 写 1 清除 \hyperref[sec:mcpwm-interrupts]{MCPWM\_TZ2\_CBC\_INT}。 (WT)
\label{fielddesc:MCPWMTZ0OSTINTCLR}\item [MCPWM\_TZ0\_OST\_INT\_CLR] 写 1 清除 \hyperref[sec:mcpwm-interrupts]{MCPWM\_TZ0\_OST\_INT}。 (WT)
\label{fielddesc:MCPWMTZ1OSTINTCLR}\item [MCPWM\_TZ1\_OST\_INT\_CLR] 写 1 清除 \hyperref[sec:mcpwm-interrupts]{MCPWM\_TZ1\_OST\_INT}。 (WT)
\label{fielddesc:MCPWMTZ2OSTINTCLR}\item [MCPWM\_TZ2\_OST\_INT\_CLR] 写 1 清除 \hyperref[sec:mcpwm-interrupts]{MCPWM\_TZ2\_OST\_INT}。 (WT)
\label{fielddesc:MCPWMCAP0INTCLR}\item [MCPWM\_CAP0\_INT\_CLR] 写 1 清除 \hyperref[sec:mcpwm-interrupts]{MCPWM\_CAP0\_INT}。 (WT)
\label{fielddesc:MCPWMCAP1INTCLR}\item [MCPWM\_CAP1\_INT\_CLR] 写 1 清除 \hyperref[sec:mcpwm-interrupts]{MCPWM\_CAP1\_INT}。 (WT)
\label{fielddesc:MCPWMCAP2INTCLR}\item [MCPWM\_CAP2\_INT\_CLR] 写 1 清除 \hyperref[sec:mcpwm-interrupts]{MCPWM\_CAP2\_INT}。 (WT)
\end{reglist}\end{regdesc}
\end{register}


\begin{register}{H}{MCPWM\_EVT\_EN\_REG}{0x0120}\label{regdesc:MCPWMEVTENREG}
\regfield{(reserved)}{2}{30}{00}%
\regfield{MCPWM\_EVT\_CAP2\_EN}{1}{29}{{0}}%
\regfield{MCPWM\_EVT\_CAP1\_EN}{1}{28}{{0}}%
\regfield{MCPWM\_EVT\_CAP0\_EN}{1}{27}{{0}}%
\regfield{MCPWM\_EVT\_TZ2\_OST\_EN}{1}{26}{{0}}%
\regfield{MCPWM\_EVT\_TZ1\_OST\_EN}{1}{25}{{0}}%
\regfield{MCPWM\_EVT\_TZ0\_OST\_EN}{1}{24}{{0}}%
\regfield{MCPWM\_EVT\_TZ2\_CBC\_EN}{1}{23}{{0}}%
\regfield{MCPWM\_EVT\_TZ1\_CBC\_EN}{1}{22}{{0}}%
\regfield{MCPWM\_EVT\_TZ0\_CBC\_EN}{1}{21}{{0}}%
\regfield{MCPWM\_EVT\_F2\_CLR\_EN}{1}{20}{{0}}%
\regfield{MCPWM\_EVT\_F1\_CLR\_EN}{1}{19}{{0}}%
\regfield{MCPWM\_EVT\_F0\_CLR\_EN}{1}{18}{{0}}%
\regfield{MCPWM\_EVT\_F2\_EN}{1}{17}{{0}}%
\regfield{MCPWM\_EVT\_F1\_EN}{1}{16}{{0}}%
\regfield{MCPWM\_EVT\_F0\_EN}{1}{15}{{0}}%
\regfield{MCPWM\_EVT\_OP2\_TEB\_EN}{1}{14}{{0}}%
\regfield{MCPWM\_EVT\_OP1\_TEB\_EN}{1}{13}{{0}}%
\regfield{MCPWM\_EVT\_OP0\_TEB\_EN}{1}{12}{{0}}%
\regfield{MCPWM\_EVT\_OP2\_TEA\_EN}{1}{11}{{0}}%
\regfield{MCPWM\_EVT\_OP1\_TEA\_EN}{1}{10}{{0}}%
\regfield{MCPWM\_EVT\_OP0\_TEA\_EN}{1}{9}{{0}}%
\regfield{MCPWM\_EVT\_TIMER2\_TEP\_EN}{1}{8}{{0}}%
\regfield{MCPWM\_EVT\_TIMER1\_TEP\_EN}{1}{7}{{0}}%
\regfield{MCPWM\_EVT\_TIMER0\_TEP\_EN}{1}{6}{{0}}%
\regfield{MCPWM\_EVT\_TIMER2\_TEZ\_EN}{1}{5}{{0}}%
\regfield{MCPWM\_EVT\_TIMER1\_TEZ\_EN}{1}{4}{{0}}%
\regfield{MCPWM\_EVT\_TIMER0\_TEZ\_EN}{1}{3}{{0}}%
\regfield{MCPWM\_EVT\_TIMER2\_STOP\_EN}{1}{2}{{0}}%
\regfield{MCPWM\_EVT\_TIMER1\_STOP\_EN}{1}{1}{{0}}%
\regfield{MCPWM\_EVT\_TIMER0\_STOP\_EN}{1}{0}{{0}}%
\reglabel{Reset}\regnewline%
\begin{regdesc}\begin{reglist}
\label{fielddesc:MCPWMEVTTIMER0STOPEN}\item [MCPWM\_EVT\_TIMER0\_STOP\_EN] 配置是否使能当定时器 0 停止时的事件生成。\\0：关闭 \\1：使能 \\(R/W)
\label{fielddesc:MCPWMEVTTIMER1STOPEN}\item [MCPWM\_EVT\_TIMER1\_STOP\_EN] 配置是否使能当定时器 1 停止时的事件生成。\\0：关闭 \\1：使能 \\(R/W)
\label{fielddesc:MCPWMEVTTIMER2STOPEN}\item [MCPWM\_EVT\_TIMER2\_STOP\_EN] 配置是否使能当定时器 2 停止时的事件生成。\\0：关闭 \\1：使能 \\(R/W)
\label{fielddesc:MCPWMEVTTIMER0TEZEN}\item [MCPWM\_EVT\_TIMER0\_TEZ\_EN] 配置是否使能当定时器 0 等于零时的事件生成。\\0：关闭 \\1：使能 \\(R/W)
\label{fielddesc:MCPWMEVTTIMER1TEZEN}\item [MCPWM\_EVT\_TIMER1\_TEZ\_EN] 配置是否使能当定时器 1 等于零时的事件生成。\\0：关闭 \\1：使能 \\(R/W)
\label{fielddesc:MCPWMEVTTIMER2TEZEN}\item [MCPWM\_EVT\_TIMER2\_TEZ\_EN] 配置是否使能当定时器 2 等于零时的事件生成。\\0：关闭 \\1：使能 \\(R/W)
\label{fielddesc:MCPWMEVTTIMER0TEPEN}\item [MCPWM\_EVT\_TIMER0\_TEP\_EN] 配置是否使能当定时器 0 等于周期时的事件生成。\\0：关闭 \\1：使能 \\(R/W)
\label{fielddesc:MCPWMEVTTIMER1TEPEN}\item [MCPWM\_EVT\_TIMER1\_TEP\_EN] 配置是否使能当定时器 1 等于周期时的事件生成。\\0：关闭 \\1：使能 \\(R/W)
\item [见下页...]
\end{reglist}\end{regdesc}
\end{register}
\addtocounter{Regfloat}{-1}
\begin{register}{H}{MCPWM\_EVT\_EN\_REG}{0x0120}
\begin{regdesc}\begin{reglist}
\item [...接上页]
\label{fielddesc:MCPWMEVTTIMER2TEPEN}\item [MCPWM\_EVT\_TIMER2\_TEP\_EN] 配置是否使能当定时器 2 等于周期时的事件生成。\\0：关闭 \\1：使能 \\(R/W)
\label{fielddesc:MCPWMEVTOP0TEAEN}\item [MCPWM\_EVT\_OP0\_TEA\_EN] 配置是否使能当生成器 0 等于时间戳 A 时的事件生成。\\0：关闭 \\1：使能 \\(R/W)
\label{fielddesc:MCPWMEVTOP1TEAEN}\item [MCPWM\_EVT\_OP1\_TEA\_EN] 配置是否使能当生成器 1 等于时间戳 A 时的事件生成。\\0：关闭 \\1：使能 \\(R/W)
\label{fielddesc:MCPWMEVTOP2TEAEN}\item [MCPWM\_EVT\_OP2\_TEA\_EN] 配置是否使能当生成器 2 等于时间戳 A 时的事件生成。\\0：关闭 \\1：使能 \\(R/W)
\label{fielddesc:MCPWMEVTOP0TEBEN}\item [MCPWM\_EVT\_OP0\_TEB\_EN] 配置是否使能当生成器 0 等于时间戳 B 时的事件生成。\\0：关闭 \\1：使能 \\(R/W)
\label{fielddesc:MCPWMEVTOP1TEBEN}\item [MCPWM\_EVT\_OP1\_TEB\_EN] 配置是否使能当生成器 1 等于时间戳 B 时的事件生成。\\0：关闭 \\1：使能 \\(R/W)
\label{fielddesc:MCPWMEVTOP2TEBEN}\item [MCPWM\_EVT\_OP2\_TEB\_EN] 配置是否使能当生成器 2 等于时间戳 B 时的事件生成。\\0：关闭 \\1：使能 \\(R/W)
\label{fielddesc:MCPWMEVTF0EN}\item [MCPWM\_EVT\_F0\_EN] 配置是否使能 FAULT0 事件生成。\\0：关闭 \\1：使能 \\(R/W)
\label{fielddesc:MCPWMEVTF1EN}\item [MCPWM\_EVT\_F1\_EN] 配置是否使能 FAULT1 事件生成。\\0：关闭 \\1：使能 \\(R/W)
\item [见下页...]
\end{reglist}\end{regdesc}
\end{register}
\addtocounter{Regfloat}{-1}
\begin{register}{H}{MCPWM\_EVT\_EN\_REG}{0x0120}
\begin{regdesc}\begin{reglist}
\item [...接上页]
\label{fielddesc:MCPWMEVTF2EN}\item [MCPWM\_EVT\_F2\_EN] 配置是否使能 FAULT2 事件生成。\\0：关闭 \\1：使能 \\(R/W)
\label{fielddesc:MCPWMEVTF0CLREN}\item [MCPWM\_EVT\_F0\_CLR\_EN] 配置是否使能 FAULT0 清除事件生成。\\0：关闭 \\1：使能 \\(R/W)
\label{fielddesc:MCPWMEVTF1CLREN}\item [MCPWM\_EVT\_F1\_CLR\_EN] 配置是否使能 FAULT0 清除事件生成。\\0：关闭 \\1：使能 \\(R/W)
\label{fielddesc:MCPWMEVTF2CLREN}\item [MCPWM\_EVT\_F2\_CLR\_EN] 配置是否使能 FAULT0 清除事件生成。\\0：关闭 \\1：使能 \\(R/W)
\label{fielddesc:MCPWMEVTTZ0CBCEN}\item [MCPWM\_EVT\_TZ0\_CBC\_EN] 配置是否使能 CBC0 事件生成。\\0：关闭 \\1：使能 \\(R/W)
\label{fielddesc:MCPWMEVTTZ1CBCEN}\item [MCPWM\_EVT\_TZ1\_CBC\_EN] 配置是否使能 CBC1 事件生成。\\0：关闭 \\1：使能 \\(R/W)
\label{fielddesc:MCPWMEVTTZ2CBCEN}\item [MCPWM\_EVT\_TZ2\_CBC\_EN] 配置是否使能 CBC2 事件生成。\\0：关闭 \\1：使能 \\(R/W)
\label{fielddesc:MCPWMEVTTZ0OSTEN}\item [MCPWM\_EVT\_TZ0\_OST\_EN] 配置是否使能 OST0 事件生成。\\0：关闭 \\1：使能 \\(R/W)
\label{fielddesc:MCPWMEVTTZ1OSTEN}\item [MCPWM\_EVT\_TZ1\_OST\_EN] 配置是否使能 OST1 事件生成。\\0：关闭 \\1：使能 \\(R/W)
\item [见下页...]
\end{reglist}\end{regdesc}
\end{register}
\addtocounter{Regfloat}{-1}
\begin{register}{H}{MCPWM\_EVT\_EN\_REG}{0x0120}
\begin{regdesc}\begin{reglist}
\item [...接上页]
\label{fielddesc:MCPWMEVTTZ2OSTEN}\item [MCPWM\_EVT\_TZ2\_OST\_EN] 配置是否使能 OST2 事件生成。\\0：关闭 \\1：使能 \\(R/W)
\label{fielddesc:MCPWMEVTCAP0EN}\item [MCPWM\_EVT\_CAP0\_EN] 配置是否使能 capture0 事件生成。\\0：关闭 \\1：使能 \\(R/W)
\label{fielddesc:MCPWMEVTCAP1EN}\item [MCPWM\_EVT\_CAP1\_EN] 配置是否使能 capture1 事件生成。\\0：关闭 \\1：使能 \\(R/W)
\label{fielddesc:MCPWMEVTCAP2EN}\item [MCPWM\_EVT\_CAP2\_EN] 配置是否使能 capture2 事件生成。\\0：关闭 \\1：使能 \\(R/W)
\end{reglist}\end{regdesc}
\end{register}


\begin{register}{H}{MCPWM\_EVT\_EN2\_REG}{0x0128}\label{regdesc:MCPWMEVTEN2REG}
\regfield{(reserved)}{26}{6}{00000000000000000000000000}%
\regfield{MCPWM\_EVT\_OP2\_TEE2\_EN}{1}{5}{{0}}%
\regfield{MCPWM\_EVT\_OP1\_TEE2\_EN}{1}{4}{{0}}%
\regfield{MCPWM\_EVT\_OP0\_TEE2\_EN}{1}{3}{{0}}%
\regfield{MCPWM\_EVT\_OP2\_TEE1\_EN}{1}{2}{{0}}%
\regfield{MCPWM\_EVT\_OP1\_TEE1\_EN}{1}{1}{{0}}%
\regfield{MCPWM\_EVT\_OP0\_TEE1\_EN}{1}{0}{{0}}%
\reglabel{Reset}\regnewline%
\begin{regdesc}\begin{reglist}
\label{fielddesc:MCPWMEVTOPNTEE1EN}\item [MCPWM\_EVT\_OP\regindex{n}\_TEE1\_EN] 配置是否在 PWM 生成器计数器的值等于 MCPWM\_OP\regindex{n}\_TSTMP\_E1\_REG 寄存器的值时，生成 MCPWM\_EVT\_OP\regindex{n}\_TEE1 事件。
\\0：不生成 \\1：生成 \\(R/W)
\label{fielddesc:MCPWMEVTOPNTEE2EN}\item [MCPWM\_EVT\_OP\regindex{n}\_TEE2\_EN] 配置是否在 PWM 生成器计数器的值等于 MCPWM\_OP\regindex{n}\_TSTMP\_E2\_REG 寄存器的值时，生成 MCPWM\_EVT\_OP\regindex{n}\_TEE2 事件。\\0：不生成 \\1：生成 \\(R/W)
\end{reglist}\end{regdesc}
\end{register}


\begin{register}{H}{MCPWM\_TASK\_EN\_REG}{0x0124}\label{regdesc:MCPWMTASKENREG}
\regfield{(reserved)}{10}{22}{0000000000}%
\regfield{MCPWM\_TASK\_CAP2\_EN}{1}{21}{{0}}%
\regfield{MCPWM\_TASK\_CAP1\_EN}{1}{20}{{0}}%
\regfield{MCPWM\_TASK\_CAP0\_EN}{1}{19}{{0}}%
\regfield{MCPWM\_TASK\_CLR2\_OST\_EN}{1}{18}{{0}}%
\regfield{MCPWM\_TASK\_CLR1\_OST\_EN}{1}{17}{{0}}%
\regfield{MCPWM\_TASK\_CLR0\_OST\_EN}{1}{16}{{0}}%
\regfield{MCPWM\_TASK\_TZ2\_OST\_EN}{1}{15}{{0}}%
\regfield{MCPWM\_TASK\_TZ1\_OST\_EN}{1}{14}{{0}}%
\regfield{MCPWM\_TASK\_TZ0\_OST\_EN}{1}{13}{{0}}%
\regfield{MCPWM\_TASK\_TIMER2\_PERIOD\_UP\_EN}{1}{12}{{0}}%
\regfield{MCPWM\_TASK\_TIMER1\_PERIOD\_UP\_EN}{1}{11}{{0}}%
\regfield{MCPWM\_TASK\_TIMER0\_PERIOD\_UP\_EN}{1}{10}{{0}}%
\regfield{MCPWM\_TASK\_TIMER2\_SYNC\_EN}{1}{9}{{0}}%
\regfield{MCPWM\_TASK\_TIMER1\_SYNC\_EN}{1}{8}{{0}}%
\regfield{MCPWM\_TASK\_TIMER0\_SYNC\_EN}{1}{7}{{0}}%
\regfield{MCPWM\_TASK\_GEN\_STOP\_EN}{1}{6}{{0}}%
\regfield{MCPWM\_TASK\_CMPR2\_B\_UP\_EN}{1}{5}{{0}}%
\regfield{MCPWM\_TASK\_CMPR1\_B\_UP\_EN}{1}{4}{{0}}%
\regfield{MCPWM\_TASK\_CMPR0\_B\_UP\_EN}{1}{3}{{0}}%
\regfield{MCPWM\_TASK\_CMPR2\_A\_UP\_EN}{1}{2}{{0}}%
\regfield{MCPWM\_TASK\_CMPR1\_A\_UP\_EN}{1}{1}{{0}}%
\regfield{MCPWM\_TASK\_CMPR0\_A\_UP\_EN}{1}{0}{{0}}%
\reglabel{Reset}\regnewline%
\begin{regdesc}\begin{reglist}
\label{fielddesc:MCPWMTASKCMPR0AUPEN}\item [MCPWM\_TASK\_CMPR0\_A\_UP\_EN] 配置是否接收 PWM 生成器 0 时间戳 A 的影子寄存器的更新任务。 \\0：无效\\1：接收\\ (R/W)
\label{fielddesc:MCPWMTASKCMPR1AUPEN}\item [MCPWM\_TASK\_CMPR1\_A\_UP\_EN] 配置是否接收 PWM 生成器 1 时间戳 A 的影子寄存器的更新任务。 \\0：无效\\1：接收\\ (R/W)
\label{fielddesc:MCPWMTASKCMPR2AUPEN}\item [MCPWM\_TASK\_CMPR2\_A\_UP\_EN] 配置是否接收 PWM 生成器 2 时间戳 A 的影子寄存器的更新任务。 \\0：无效\\1：接收\\ (R/W)
\label{fielddesc:MCPWMTASKCMPR0BUPEN}\item [MCPWM\_TASK\_CMPR0\_B\_UP\_EN] 配置是否接收 PWM 生成器 0 时间戳 B 的影子寄存器的更新任务。 \\0：无效\\1：接收\\ (R/W)
\label{fielddesc:MCPWMTASKCMPR1BUPEN}\item [MCPWM\_TASK\_CMPR1\_B\_UP\_EN] 配置是否接收 PWM 生成器 1 时间戳 B 的影子寄存器的更新任务。 \\0：无效\\1：接收\\ (R/W)
\label{fielddesc:MCPWMTASKCMPR2BUPEN}\item [MCPWM\_TASK\_CMPR2\_B\_UP\_EN] 配置是否接收 PWM 生成器 2 时间戳 B 的影子寄存器的更新任务。 \\0：无效\\1：接收\\ (R/W)
\item [见下页...]
\end{reglist}\end{regdesc}
\end{register}
\addtocounter{Regfloat}{-1}
\begin{register}{H}{MCPWM\_TASK\_EN\_REG}{0x0124}
\begin{regdesc}\begin{reglist}
\item [...接上页]
\label{fielddesc:MCPWMTASKGENSTOPEN}\item [MCPWM\_TASK\_GEN\_STOP\_EN] 配置是否接收所有 PWM 生成器的停止任务。 \\0：无效\\1：接收\\ (R/W)
\label{fielddesc:MCPWMTASKTIMER0SYNCEN}\item [MCPWM\_TASK\_TIMER0\_SYNC\_EN] 配置是否接收定时器 0 的同步任务。 \\0：无效\\1：接收\\ (R/W)
\label{fielddesc:MCPWMTASKTIMER1SYNCEN}\item [MCPWM\_TASK\_TIMER1\_SYNC\_EN] 配置是否接收定时器 1 的同步任务。 \\0：无效\\1：接收\\ (R/W)
\label{fielddesc:MCPWMTASKTIMER2SYNCEN}\item [MCPWM\_TASK\_TIMER2\_SYNC\_EN] 配置是否接收定时器 2 的同步任务。 \\0：无效\\1：接收\\ (R/W)
\label{fielddesc:MCPWMTASKTIMER0PERIODUPEN}\item [MCPWM\_TASK\_TIMER0\_PERIOD\_UP\_EN] 配置是否接收定时器 0 的周期更新任务。 \\0：无效\\1：接收\\ (R/W)
\label{fielddesc:MCPWMTASKTIMER1PERIODUPEN}\item [MCPWM\_TASK\_TIMER1\_PERIOD\_UP\_EN] 配置是否接收定时器 1 的周期更新任务。 \\0：无效\\1：接收\\ (R/W)
\label{fielddesc:MCPWMTASKTIMER2PERIODUPEN}\item [MCPWM\_TASK\_TIMER2\_PERIOD\_UP\_EN] 配置是否接收定时器 2 的周期更新任务。 \\0：无效\\1：接收\\ (R/W)
\label{fielddesc:MCPWMTASKTZ0OSTEN}\item [MCPWM\_TASK\_TZ0\_OST\_EN] 配置是否接收 OST0 任务。 \\0：无效\\1：接收\\ (R/W)
\label{fielddesc:MCPWMTASKTZ1OSTEN}\item [MCPWM\_TASK\_TZ1\_OST\_EN] 配置是否接收 OST1 任务。 \\0：无效\\1：接收\\ (R/W)
\item [见下页...]
\end{reglist}\end{regdesc}
\end{register}
\addtocounter{Regfloat}{-1}
\begin{register}{H}{MCPWM\_TASK\_EN\_REG}{0x0124}
\begin{regdesc}\begin{reglist}
\item [...接上页]
\label{fielddesc:MCPWMTASKTZ2OSTEN}\item [MCPWM\_TASK\_TZ2\_OST\_EN] 配置是否接收 OST2 任务。 \\0：无效\\1：接收\\ (R/W)
\label{fielddesc:MCPWMTASKCLR0OSTEN}\item [MCPWM\_TASK\_CLR0\_OST\_EN] 配置是否接收 OST0 清除任务。 \\0：无效\\1：接收\\ (R/W)
\label{fielddesc:MCPWMTASKCLR1OSTEN}\item [MCPWM\_TASK\_CLR1\_OST\_EN] 配置是否接收 OST1 清除任务。 \\0：无效\\1：接收\\ (R/W)
\label{fielddesc:MCPWMTASKCLR2OSTEN}\item [MCPWM\_TASK\_CLR2\_OST\_EN] 配置是否接收 OST2 清除任务。 \\0：无效\\1：接收\\ (R/W)
\label{fielddesc:MCPWMTASKCAP0EN}\item [MCPWM\_TASK\_CAP0\_EN] 配置是否接收 capture0 任务。 \\0：无效\\1：接收\\ (R/W)
\label{fielddesc:MCPWMTASKCAP1EN}\item [MCPWM\_TASK\_CAP1\_EN] 配置是否接收 capture1 任务。 \\0：无效\\1：接收\\ (R/W)
\label{fielddesc:MCPWMTASKCAP2EN}\item [MCPWM\_TASK\_CAP2\_EN] 配置是否接收 capture2 任务。 \\0：无效\\1：接收\\ (R/W)
\end{reglist}\end{regdesc}
\end{register}


\begin{register}{H}{MCPWM\_OP\regindex{n}\_TSTMP\_E1\_REG (\regindex{n}: 0-2)}{0x012C+0x8*\regindex{n}}\label{regdesc:MCPWMOPNTSTMPE1REG}
\regfield{(reserved)}{16}{16}{0000000000000000}%
\regfield{MCPWM\_OP\regindex{n}\_TSTMP\_E1}{16}{0}{{0}}%
\reglabel{Reset}\regnewline%
\begin{regdesc}\begin{reglist}
\label{fielddesc:MCPWMOPNTSTMPE1}\item [MCPWM\_OP\regindex{n}\_TSTMP\_E1] 配置生成器\regindex{n} 时间戳 E1 值。 (R/W)
\end{reglist}\end{regdesc}
\end{register}


\begin{register}{H}{MCPWM\_OP\regindex{n}\_TSTMP\_E2\_REG(\regindex{n}: 0-2)}
{0x0130+0x8*\regindex{n}}\label{regdesc:MCPWMOPNTSTMPE2REG}
\regfield{(reserved)}{16}{16}{0000000000000000}%
\regfield{MCPWM\_OP\regindex{n}\_TSTMP\_E2}{16}{0}{{0}}%
\reglabel{Reset}\regnewline%
\begin{regdesc}\begin{reglist}
\label{fielddesc:MCPWMOPNTSTMPE2}\item [MCPWM\_OP\regindex{n}\_TSTMP\_E2] 配置生成器\regindex{n} 时间戳 E2 值。(R/W)
\end{reglist}\end{regdesc}
\end{register}


\begin{register}{H}{MCPWM\_CLK\_REG}{0x0144}\label{regdesc:MCPWMCLKREG}
\regfield{(reserved)}{31}{1}{0000000000000000000000000000000}%
\regfield{MCPWM\_CLK\_EN}{1}{0}{{0}}%
\reglabel{Reset}\regnewline%
\begin{regdesc}\begin{reglist}
\label{fielddesc:MCPWMCLKEN}\item [MCPWM\_CLK\_EN] 配置是否强制使能时钟。\\
0：无效 \\
1：强制使能时钟 \\
(R/W)
\end{reglist}\end{regdesc}
\end{register}


\begin{register}{H}{MCPWM\_VERSION\_REG}{0x0148}\label{regdesc:MCPWMVERSIONREG}
\regfield{(reserved)}{4}{28}{0000}%
\regfield{MCPWM\_DATE}{28}{0}{{0x2107230}}%
\reglabel{Reset}\regnewline%
\begin{regdesc}\begin{reglist}
\label{fielddesc:MCPWMDATE}\item [MCPWM\_DATE] 版本控制寄存器。 (R/W)
\end{reglist}\end{regdesc}
\end{register}


